\documentclass[11pt,letterpaper]{article}
\usepackage[T1]{fontenc}
\usepackage{lmodern,microtype}
\usepackage[margin=1in]{geometry}
\usepackage{amsmath,amssymb,amsthm,mathtools,booktabs,array,longtable}
\usepackage{xcolor,enumitem,fancyhdr}
\usepackage[colorlinks=true,linkcolor=blue!45!black,citecolor=blue!45!black,urlcolor=blue!45!black,breaklinks=true]{hyperref}
\pdfmapfile{+lm.map}\pdfmapfile{+cm.map}\pdfmapfile{+symbols.map}
\pdfinfoomitdate=1\pdfsuppressptexinfo=15
\pdftrailerid{<53796D5061636B65643230323631303035><53796D5061636B65643230323631303035>}
\providecommand{\ReportNumber}{238}
\hypersetup{pdftitle={All Fixed Order Asymptotics for Symmetric Packed Matrices},pdfauthor={Report \ReportNumber},pdfsubject={OEIS A138178 marked expansions collision laws and controlled inverses}}
\pagestyle{fancy}\fancyhf{}\fancyhead[L]{\small Symmetric packed matrices}\fancyhead[R]{\small Report \ReportNumber}\fancyfoot[C]{\thepage}
\renewcommand{\headrulewidth}{0pt}\setlength{\headheight}{14pt}
\setlength{\parskip}{4pt}\setlength{\emergencystretch}{2em}
\setlist{itemsep=2pt,topsep=4pt}\numberwithin{equation}{section}
\newtheorem{theorem}{Theorem}[section]\newtheorem{lemma}[theorem]{Lemma}\newtheorem{proposition}[theorem]{Proposition}\newtheorem{corollary}[theorem]{Corollary}
\theoremstyle{definition}\newtheorem{definition}[theorem]{Definition}
\theoremstyle{remark}\newtheorem{remark}[theorem]{Remark}
\newcommand{\E}{\mathbb E}\newcommand{\PP}{\mathbb P}\newcommand{\R}{\mathbb R}\newcommand{\C}{\mathbb C}\newcommand{\N}{\mathbb N}\newcommand{\ii}{\mathrm i}\newcommand{\dd}{\mathrm d}\newcommand{\eps}{\varepsilon}\newcommand{\ceil}[1]{\left\lceil #1\right\rceil}\newcommand{\TV}{d_{\mathrm{TV}}}\newcommand{\Pois}{\operatorname{Pois}}\newcommand{\Var}{\operatorname{Var}}\newcommand{\Cov}{\operatorname{Cov}}\newcommand{\Law}{\operatorname{Law}}
\begin{document}
\begin{center}
{\large Report \ReportNumber}\par\vspace{8pt}
{\LARGE\bfseries All Fixed Order Asymptotics\\[5pt] for Symmetric Packed Matrices}\par\vspace{9pt}
{\large Dimension, trace, collision laws, and controlled inverses}\par\vspace{7pt}
5 October 2026
\end{center}
\begin{abstract}
Let $A_n$ count symmetric nonnegative integer matrices of total entry sum $n$ with no zero line and with ordered, distinguished positions, the sequence OEIS A138178. Writing $L=\log 2$ and $I_n$ for the involution number, we prove
\[
 A_n=\frac{e^{L^2/4}I_n}{2L^{n+1}}
 \left(1-\frac{L^2}{6\sqrt n}
 +\frac{-L^2/4+L^4/18}{n}+O(n^{-3/2})\right),
\]
and construct every fixed subsequent order. The proof is uniform for complex dimension and trace markers near one. It controls all polynomial degrees in the geometric transform, the full auxiliary-dimension range, and both angular saddles before expanding a normalized pair of Gaussian integrals. Dimension and trace have independent Gaussian limits on the respective scales $n^{1/2}$ and $n^{1/4}$. A further theorem, uniform on every fixed compact set of collision markers including zero, gives independent Poisson counts of repeated diagonal and off-diagonal cells, their joint independence from the Gaussian pair, a sharp leading total-variation constant, and an $O(n^{-1})$ corrected Poisson approximation. We also give an all-fixed-order Lambert-$W$ inverse with rounding-safe two-ceiling bounds. The geometric transform is attributed to Jovovic in OEIS; the leading symmetric binary equivalent and a related nonsymmetric collision mechanism are prior results of Cameron, Prellberg, and Stark. No global priority, complete transseries, or effective numerical onset is asserted.
\end{abstract}
\setcounter{tocdepth}{1}
\begingroup\small\setlength{\parskip}{0pt}
\makeatletter\renewcommand\l@section[2]{\@dottedtocline{1}{0em}{2em}{#1}{#2}}\makeatother
\tableofcontents\endgroup
\clearpage
\section{The model and its exact transform}\label{spm:sec:model}
A matrix is \emph{packed} if it has no zero row or column. Let $\mathcal A_n$ be the set of symmetric packed matrices with nonnegative integer entries and total entry sum $n$. The dimension is allowed to vary. Rows and columns are ordered, and simultaneous permutations are not identified. The empty matrix is the sole object of weight zero. For $n>0$, the dimension is at most $n$, so the set is finite. An off-diagonal entry of value $m$ contributes $2m$ to the total sum, while a diagonal entry contributes its value once.

Put $A_n=|\mathcal A_n|$. The first values are
\[
1,\ 1,\ 3,\ 9,\ 33,\ 125,\ 531,\ 2349,\ 11205,\ 55589,\ 291423.
\]
This is OEIS A138178~\cite{oeis138178}. The entry also identifies these objects with normal semistandard Young tableaux, whose used alphabet is an initial interval. The matrix model, rather than any quotient by relabeling, is used throughout this report.

For $M\in\mathcal A_n$, let $D(M)$ be its actual dimension and $T(M)$ its diagonal sum, and define
\[
 A_n(u,v)=\sum_{M\in\mathcal A_n}u^{D(M)}v^{T(M)},\qquad A_0(u,v)=1.
\]
The auxiliary quantities
\begin{equation}\label{spm:eq:notation}
 q=\frac{u}{1+u},\qquad L=L(u)=\log(1+1/u),\qquad
 I_n(v)=n![x^n]e^{vx+x^2/2}
\end{equation}
use the branch of the logarithm positive at $u=1$. In particular $L(1)=\log2$, and $I_n(1)$ is the involution number. Define the coefficient polynomial
\begin{equation}\label{spm:eq:fn}
 f_n(k,v)=[z^n](1-vz)^{-k}(1-z^2)^{-(k^2-k)/2}.
\end{equation}
For integer $k\geq0$, this counts all symmetric $k\times k$ matrices, with zero lines allowed, marking their diagonal sum. For every fixed $n>0$, it is a polynomial in $k$ of degree at most $n$ with zero constant term.

\begin{proposition}[Exact geometric transform]\label{spm:prop:transform}
For $u,v$ in a complex neighborhood of $(1,1)$,
\begin{equation}\label{spm:eq:transform}
 A_n(u,v)=\frac1{1+u}\sum_{k\geq0}q^k f_n(k,v).
\end{equation}
The series converges absolutely for fixed $n$. The marker $u$ records the actual packed dimension, not the auxiliary index $k$.
\end{proposition}
\begin{proof}
Deleting simultaneous zero rows and columns from a $k\times k$ matrix leaves a unique packed matrix. Each particular packed matrix of dimension $d$ has $\binom{k}{d}$ preimages. Since $1-q=(1+u)^{-1}$,
\[
 (1-q)\sum_{k\geq d}q^k\binom{k}{d}
 =\left(\frac{q}{1-q}\right)^d=u^d.
\]
The diagonal sum is unchanged by insertion of zero lines. Summing over the finite family of packed matrices proves the identity. The polynomial degree assertion follows by expanding the logarithm of \eqref{spm:eq:fn}: every atom $k^a z^b$ has $a\leq b$.
\end{proof}

The unmarked geometric generating-function transform is already recorded in OEIS, credited to Vladeta Jovovic in December 2009~\cite{oeis138178}. Its use here, and the short marked extension above, must not be mistaken for a newly discovered unmarked identity.

\subsection{Prior results and the boundary of the model}
Cameron, Prellberg, and Stark~\cite{cps}, Proposition 4.2, prove the leading equivalent for the symmetric \emph{binary} companion $B_n$, OEIS A135588~\cite{oeis135588}:
\begin{equation}\label{spm:eq:priorbinary}
 B_n\sim e^{-L-L^2/4}\frac{I_n(1)P_n}{n!},\qquad L=\log2,
\end{equation}
where $P_n$ is the ordered Bell number. Their Section 2 also obtains a Poisson collision limit in a nonsymmetric incidence-matrix argument based on pairs of random preorders. The collision mechanism therefore has substantial precedent. The present treatment concerns the nonnegative symmetric model, uniform marked expansions to every fixed order, and quantitative refinements of its joint collision law.

Rectangular packed matrices, matrices with decreasing margins, and matrices considered up to row or column permutations are different counting models. Related rectangular and Burge-polynomial frameworks appear in \cite{matrixcomp,caylerian}. Their familiar names and related generating functions do not identify their sequences with $A_n$. A bounded check of relevant public sources and existing reports found no duplicate of the result assembled here; it is not a certificate of worldwide novelty. The source notes accompanying this report describe the public attributions and the limits of that check.

\begin{writenote}[Added 6 October 2026, batch~108: the rectangular counterpart]
The rectangular counterpart of this model on the diagonal of total weight
$2N$ with $N$ rows (OEIS A261784; bundle Report~169, which the source notes
name among the reports they read) is Part~IV (Sections~34--44, labels
\file{mxc:pm:}) of the collection's report
\file{a261781-matrix-compositions}. It uses the devices of this report on a
different model: a geometric weighting of dummy columns that erases zero
lines, the pole of the ordered Bell generating function at $\log(1+1/u)$,
an expansion to every fixed order uniform for complex markers of the
columns and of the cells with entry at least~$2$ (its
Theorem~\file{mxc:pm:thm:marked}, on $|\omega-1|\le2.1$, $|u-1|\le0.02$),
and a joint limit in which the number of such cells is Poisson of mean
$r\log2=1.1046\ldots$, $r=2+W_0(-2e^{-2})$, independent of the Gaussian
column count (Theorem~\file{mxc:pm:thm:joint}). There the expansion is in
powers of $N^{-1}$; here it is in powers of $n^{-1/2}$ around an
involution saddle, and the repeated cells split into diagonal and
off-diagonal ones with means $\log2$ and $(\log2)^2/2$
(Theorem~\ref{spm:thm:joint}). No theorem is shared, and neither report
uses the other. The collection's report
\file{a260700-parabolic-double-cosets} extracts the same ordered Bell pole
at $\log2$ for a different sequence, and the exact pole-lattice formula
for the ordered Bell numbers is Theorem~\file{q2:thm:fubini} of the
transseries volume
\path{Analysis/Transseries/docs/series-and-transseries/Transseries_And_Inversion/transseries_and_inversion.tex}:
the last sentence of Corollary~\ref{spm:cor:unmarked} is its one-pole
truncation, and~\eqref{spm:eq:poles} at $L=\log2$, $j=n$, multiplied
by~$\tfrac12$, is its first equality.
\end{writenote}

\begin{remark}[{[write]} The OEIS entries and the data, 6~October 2026]\label{spm:rem:oeis}
The live entry A138178 \cite{oeis138178} (revision~\#39 of 18~December
2022, the revision the source inspected; read 6~October 2026) is named,
verbatim, ``Number of symmetric matrices with nonnegative integer entries
and without zero rows or columns such that sum of all entries is equal to
n.'' Its offset is~$0$; its 25 data terms $A_0,\ldots,A_{24}$ equal the
source's transcription (the delivered receipt
\file{data/count_receipt.json}), the first eleven being those displayed
above. The entry has two generating functions: an unattributed
inclusion--exclusion form, which is the generating function
of~\eqref{spm:eq:deletion} at $u=v=1$, and Jovovic's
$\sum_{k\ge0}2^{-k-1}(1-x)^{-k}(1-x^2)^{-\binom k2}$ of 9~December 2009,
which is~\eqref{spm:eq:transform} at $u=v=1$ (then $q=\tfrac12$). Its
comments give Wiseman's normal semistandard tableaux (2018) and
examples; it has no asymptotic statement, so Corollary~\ref{spm:cor:unmarked}
is the only asymptotic for A138178 in the entry or in the repository.
Heinz's b-file lists $A_0,\ldots,A_{500}$. The source's retrieval failed
(\file{SOURCES.md}), and it claims no comparison
(Section~\ref{spm:sec:computation}); the write fetched it on
6~October 2026 and computed $A_0,\ldots,A_{500}$ exactly by an
implementation of its own (the recurrence~\eqref{spm:eq:exactrecurrence}
at $v=1$ and $\tfrac12\sum_{k\ge0}2^{-k}k^j=P_j$ for the ordered Bell
numbers): all 501 terms agree. The binary companion A135588
\cite{oeis135588} (revision~\#17 of 22~January 2024) has 26 data terms
$B_0,\ldots,B_{25}$; an inclusion--exclusion count of the write's
(binomial cell factors, as in Section~\ref{spm:sec:computation}) gives
all 26. It has no asymptotic statement either. Nothing in either entry is
corrected, and nothing was submitted to the OEIS. The leading formula and
its corrections are compared with the exact counts in
Remark~\ref{spm:rem:third}.

\emph{The binary precedent.} The write read Cameron--Prellberg--Stark
\cite{cps} (arXiv version~2 of 4~April 2006). Their
Proposition~4.2 (printed p.~14, proof pp.~14--15) states $S_{11}(n)\sim C_sI(n)P(n)/n!$ with
$C_s=\tfrac12e^{-(\log2)^2/4}\approx0.44341$, where $S_{11}(n)$ counts
the symmetric incidence matrices with $n$ ones (A135588, here $B_n$),
$I(n)$ the involutions and $P(n)$ the preorders, that is the ordered Bell
numbers. Since $\tfrac12=e^{-\log2}$, this is~\eqref{spm:eq:priorbinary}.
Their Section~2 (the first proof of their Theorem~2.1, pp.~6--9) shows by
the method of moments that the number $W$ of pairs of points lying in a
common block of both of two independent uniform random preorders converges
to a Poisson$((\log2)^2/2)$ law, so that $\PP(W=0)\to e^{-(\log2)^2/2}$
(their~(8), p.~9):
the nonsymmetric collision precedent the source credits. The source's
account of both is accurate.
\end{remark}

\subsection{Provenance, sources and reading conventions}\label{spm:sec:provenance}
\begin{writenote}[Added 6 October 2026, batch~108]
\emph{Provenance.} This report is Report~238 of the session bundle that
ProveIt's \file{docs/incoming} intake processed as batch~108: the archive
\file{Report238.zip} (604{,}411 bytes, SHA-256
\file{255ba5ee21d1...5f6eac3b}, 30 files in a wrapper directory
\file{Report238/}), arrival commit \file{60f54ea06}, placed in
\file{602e5bd0f}. The text is the delivered \file{article.tex} with its
twelve files \file{sections/*.tex} (952 lines in all, dated 5~October
2026), shipped under the same names. The delivered 22-page PDF and the
checksum manifest \file{MANIFEST.sha256} (29 entries, verified at
placement and again by the write) are not shipped, and the delivered
README is replaced by the report README; all three are retrievable from
the arrival commit. The package records no ProveIt commit; its
\file{SOURCES.md} describes a bounded look at existing reports, naming
the matrix-composition report and bundle Reports~169 and~177 (now Part~IV
of \file{a261781-matrix-compositions} and the batch-109 report
\file{a262810-diagonal-alignments}). The phrase ``private review'' does
not occur in it; the PDF author field is ``Report 238''. It is a
single-source report, so no merge choice was made. The write prefixed
the 110 delivered labels with \file{spm:} and updated the 93 references
to them; it added this subsection, the notes marked [write],
Remarks~\ref{spm:rem:oeis}, \ref{spm:rem:third}
and~\ref{spm:rem:transseries} and Section~\ref{spm:sec:further}. Apart
from the label prefixes no theorem, proof or number of the source was
changed: the added statements are the last of their sections, the added
section comes after the last numbered one, and the added displays are
unnumbered, so every theorem, section and equation keeps its delivered
number.

\emph{The sources, as the write read them} (6~October 2026).
\begin{itemize}\raggedright
\item OEIS A138178 (revision~\#39), its b-file, and A135588
 (revision~\#17) (Remark~\ref{spm:rem:oeis}).
\item Cameron--Prellberg--Stark \cite{cps}, arXiv:math/0510155v2:
 Proposition~4.2 and its proof (pp.~14--15) and the first proof of their
 Theorem~2.1 in Section~2 (pp.~6--9), as in Remark~\ref{spm:rem:oeis}.
\item The transseries volume
 \path{Analysis/Transseries/docs/series-and-transseries/Transseries_And_Inversion/transseries_and_inversion.tex}:
 \file{p0:def:core}, \file{p0:prop:factorial-core},
 \file{p0:thm:core-reversion} with \file{p0:rem:core-instances},
 \file{p0:def:three-inverses}, \file{p0:thm:staircase},
 \file{plt:thm:lw-template} and \file{q2:thm:fubini}, read for
 Remark~\ref{spm:rem:transseries} and the note above.
\item Part~IV of \file{a261781-matrix-compositions}, at the statements
 named in the note above.
\item Not read by the write: Munarini--Poneti--Rinaldi \cite{matrixcomp}
 and Cerbai--Claesson \cite{caylerian}. What is said about them (related
 models, not the transpose-symmetric count) is the source's.
\end{itemize}

\emph{What was checked at intake and by the write.} The batch-108 dossier
read the manuscript and found no error; it checked by hand the
total-variation constant of Theorem~\ref{spm:thm:TV} from the first mean
corrections, and reran the package on a copy (Windows, Python~3.14.4,
SymPy~1.14.0, mpmath~1.3.0): \file{build.py} with \verb|--verify-only|
verified 29 files; \file{exact_counts.py} (to $n=640$),
\file{symbolic_coefficients.py} and \file{diagnostics.py} reproduced the
three delivered receipts; \file{guard_tests.py} stops on Windows at its
own path guard; the PDF and ZIP rebuilds were not run. The write reran the
three mathematical programs from the shipped files (README, Route~B):
their outputs equal the receipts byte for byte. It re-derived by hand,
and in part with SymPy, $C_1$ and $\mathcal B_1$, $G_1$, $H_1$ of
Section~\ref{spm:sec:allorders}, the agreement of~\eqref{spm:eq:D2} with
the delivered symbolic receipt, the cumulants of
Proposition~\ref{spm:prop:cumulants} and the residue~\eqref{spm:eq:trace-residue},
the coefficients~\eqref{spm:eq:Sstar}, \eqref{spm:eq:C1star}
and~\eqref{spm:eq:binaryS}--\eqref{spm:eq:binaryprob}, the mass
correction~\eqref{spm:eq:masscorrection} and the constant
of~\eqref{spm:eq:sharpTV}, the constants~\eqref{spm:eq:cbd}, and the
coefficients~\eqref{spm:eq:alpha}--\eqref{spm:eq:gamma} (also by the
triangular recurrence of the transseries volume,
Remark~\ref{spm:rem:transseries}); and it confirmed the monotonicity
argument of Section~\ref{spm:sec:inverse}. Finite computations are
observations, not certificates.

\emph{Relation to the repository.} Before batch~108 no file of the
repository mentioned A138178, A135588 or symmetric packed matrices. No
statement of this report is formalized in Lean or Rocq, and its place in
the collection confers no formal status. The ordered Bell numbers used in
Corollary~\ref{spm:cor:unmarked} are defined in Lean as
\file{Fabius.fubini}, with their exponential generating function
(\file{Fabius.two_sub_exp_mul_egfA_fubini}), in
\path{Analysis/FabiusFunction/Lean/FabiusFunction/OrderedBell.lean}; no
asymptotic statement about them is formalized there. The neighbouring
reports are named in the note above and in the report README.
\end{writenote}
\begin{writenote}[What is not claimed, collected from the source]
The unmarked geometric transform is Jovovic's (OEIS, 2009) and ``must not
be mistaken for a newly discovered unmarked identity''; the binary leading
equivalent is Cameron--Prellberg--Stark's, ``recovered here as a
consistency check'', and the Poisson collision mechanism ``has substantial
precedent''. The expansions are local in the dimension and trace markers:
the trace marker is not carried through $v=0$, the order $J$ is fixed
independently of $n$, and ``no uniform estimate in a growing order is
asserted''. The diagnostics at $(u,v)=(2,2)$ and $(1,\tfrac12)$ are ``not
an enlargement'' of the neighbourhood of Theorem~\ref{spm:thm:main}. No
span-one local limit of $T_n$, and no total-variation approximation of
$T_n$ by an unrestricted Poisson law, is asserted. The signed corrections
are ``not claimed positive''; the adjusted Poisson law is not claimed
optimal, and neither the $O(n^{-1})$ constant nor a starting index is
effective. The inverse is ``the inverse of the specified smooth
approximation, not a definition of a noninteger combinatorial count'';
\eqref{spm:eq:rounding} ``does not claim that the ceiling of the bare
displayed truncation is always correct''; Theorem~\ref{spm:thm:ceilings}
gives no numerical constants and ``is an asymptotic bound, not a certified
finite-input inverse program''. The bounded source check ``is not a
certificate of worldwide novelty'', and no global priority, complete
transseries or effective onset is asserted. The exact checks have hard caps
that are implementation boundaries, ``not a range restriction on the
mathematical theorems''; the symbolic checks ``do not constitute a
machine-checked proof''; the decimal outputs are noncertified; and
reproducibility is claimed for the recorded toolchain only.
\end{writenote}
\medskip\noindent\begin{minipage}{\textwidth}
\emph{Reading conventions} (letters the manuscript reuses, with the
tempting false readings; no symbol of the source was renamed; symbols of
the transseries volume carry the subscript ``vol'' in
Remark~\ref{spm:rem:transseries}).\par
\smallskip
\begingroup\small\renewcommand{\arraystretch}{1.1}
\noindent\begin{tabular}{@{}>{\raggedright\arraybackslash}p{0.75in}>{\raggedright\arraybackslash}p{3.45in}>{\raggedright\arraybackslash}p{1.95in}@{}}
\toprule
Symbol & Meaning here (where) & Not to be read as\\
\midrule
$L$ & $\log(1+1/u)$, so $\log2$ at $u=1$ and throughout Sections~\ref{spm:sec:TV}--\ref{spm:sec:inverse} & the volume's target $L_{\rm vol}$\\
$q$ & $u/(1+u)=e^{-L}$ & a Poisson mean\\
$w$ & collision marker (Section~\ref{spm:sec:collision}); Gaussian angular variable $w$, $W$ (Section~\ref{spm:sec:allorders}); $W(2y/(eL^2))$ (Section~\ref{spm:sec:inverse}) & each other\\
$D(M)$, $D_n$ & dimension of a matrix, of a uniform one & $D_2$, the $n^{-1}$ log-coefficient~\eqref{spm:eq:D2}; the exponent $D_J$ in~\eqref{spm:eq:taylorbound}\\
$S(L,v)$, $S_*$; $S(M)$, $S_n$ & amplitude exponents \eqref{spm:eq:S}, \eqref{spm:eq:Sstar}; repeated off-diagonal cells & each other\\
$R$, $R_n$, $R_j$ & pole-separating radius~\eqref{spm:eq:R}; repeated diagonal cells; inverse coefficients~\eqref{spm:eq:inverseall} & each other\\
$K$ & auxiliary integer (Section~\ref{spm:sec:saddles}); compact marker set (Theorem~\ref{spm:thm:collision}) & each other; not the dimension $D$\\
$B$, $B_n$, $B_J$ & correction~\eqref{spm:eq:B}; binary count (A135588); remainder constant (Theorem~\ref{spm:thm:ceilings}) & each other\\
$a$, $b$, $c$, $d$ & in Section~\ref{spm:sec:inverse} the constants~\eqref{spm:eq:cbd}; elsewhere $a=(v^2-1)/2$ or~\eqref{spm:eq:C1star}, $c=(2\zeta-1)/4$, $a=\log u$, $b=\log v$ (Section~\ref{spm:sec:fluctuations}), $b_{nj}$, $d_m$, $d_n(k)$, and generic constants $c$, $C$ & each other; $a_{\rm vol}$, $b_{\rm vol}$, $d_{\rm vol}$\\
$h$, $H$, $H_j$, $H_\pm$ & $n^{-1/2}$; angular exponent~\eqref{spm:eq:H}; Charlier-type corrections~\eqref{spm:eq:Charlier}; envelopes in the proof of Theorem~\ref{spm:thm:ceilings} & each other; $h_{\rm vol}$\\
$Q$, $Q_j$ & $w+1$ in~\eqref{spm:eq:Wseed}; normalized collision polynomials~\eqref{spm:eq:Qpolynomials} & each other; the integer polynomials $Q_n$ of \file{COMPUTATION.md}\\
$N$ & Lambert seed~\eqref{spm:eq:Wseed}; also the bound $N\ge|\zeta-1|$ in~\eqref{spm:eq:collision-envelope} & each other; $N_\ast^{\rm vol}$\\
$M$, $M_n$ & a matrix, the bound $M\ge|w-1|$; the main term~\eqref{spm:eq:gammaexact} & each other\\
$P_n$, $P_j(u)$, $P_1$, $P_2$ & ordered Bell numbers; ordered-Bell transforms (Section~\ref{spm:sec:computation}); Poisson variables & each other\\
$\lambda$, $\mu$, $\kappa$, $\rho$ & Poisson means $L$, $L^2/2$, $\kappa=L+L^2$; $\rho(h)$ of~\eqref{spm:eq:scaling} and the radius $\rho>1$ of Section~\ref{spm:sec:TV}; $\lambda$ is also a scale in Lemma~\ref{spm:lem:gammabound} & each other; $\Lambda$ is the involution law of Section~\ref{spm:sec:saddles}, not $\Lambda_{\rm vol}$\\
$E_n$, $\mathcal E$ & pole error~\eqref{spm:eq:mainpluserror}, generating-function remainder (Section~\ref{spm:sec:TV}); Gaussian functional & each other; $\E$\\
$r$, $s$, $t$ & radius~\eqref{spm:eq:r}, Poisson arguments, gamma variable; also markers in proofs & each other\\
\bottomrule
\end{tabular}
\endgroup
\end{minipage}\par

\section{The all fixed order marked expansion}\label{spm:sec:theorem}
Write
\begin{align}
 S(L,v)&=\frac{(v^2-1)L}{2}+\frac{L^2}{4},\label{spm:eq:S}\\
 C_1(L,v)&=-\frac{v(v^2-1)L}{2}+\frac{v(2v^2-3)L^2}{6},\label{spm:eq:C1}\\
 D_2(L,v)&=\frac{v^2(v^2-1)(L+L^3)}4
          +\frac{(-5v^4+6v^2-3)L^2}{8}+\frac{L^4}{24},\label{spm:eq:D2}\\
 C_2(L,v)&=D_2(L,v)+\frac12C_1(L,v)^2.\label{spm:eq:C2}
\end{align}
Here $D_2$ is the coefficient of $n^{-1}$ in the logarithm of the normalized expansion, which makes several subsequent formulas shorter.

\begin{theorem}[Uniform marked expansion]\label{spm:thm:main}
There is a complex neighborhood $U$ of $(1,1)$ such that, for every fixed integer $J\geq0$, uniformly on compact subsets of $U$,
\begin{equation}\label{spm:eq:main}
 A_n(u,v)=\frac{I_n(v)e^{S(L,v)}}{(1+u)L^{n+1}}
 \left(\sum_{j=0}^{J}C_j(L,v)n^{-j/2}
       +O_J(n^{-(J+1)/2})\right).
\end{equation}
The coefficients $C_j$ are polynomials in $L,v$ with rational coefficients, $C_0=1$, and the first two are \eqref{spm:eq:C1}--\eqref{spm:eq:C2}. A finite Gaussian-moment prescription for every $C_j$ is given in \eqref{spm:eq:algorithm}. On a smaller compact set, every fixed derivative of the normalized holomorphic remainder obeys the same order bound.
\end{theorem}

The assertion is local in the dimension and trace markers. In particular it does not extend the trace marker through $v=0$, where an exact parity boundary and a second comparable involution saddle arise. The order $J$ is fixed independently of $n$; no uniform estimate in a growing order is asserted.

\begin{corollary}[Unmarked expansion]\label{spm:cor:unmarked}
With $L=\log2$,
\begin{equation}\label{spm:eq:unmarked}
 A_n=\frac{e^{L^2/4}I_n(1)}{2L^{n+1}}
 \left(1-\frac{L^2}{6\sqrt n}
 +\frac{-L^2/4+L^4/18}{n}+O(n^{-3/2})\right).
\end{equation}
The factor $1/(2L^{n+1})$ can be replaced by $P_n/n!$ with exponentially small relative error.
\end{corollary}
\begin{proof}
Substitute $u=v=1$ in the theorem. The ordered Bell exponential generating function is $(2-e^z)^{-1}$; its poles are $L+2\pi\ii m$. The pole at $L$ gives $P_n/n!=1/(2L^{n+1})$ up to an exponentially small relative error. Equivalently, apply the exact moment formula in \eqref{spm:eq:poles} to $\frac12\sum_{k\geq0}2^{-k}k^n$.
\end{proof}

The proof has three necessary parts. First an everywhere-positive coefficient envelope controls the entire nonreal-pole error, including low polynomial degrees. Next a gamma variable and a complex involution contour are localized with absolute estimates over their full ranges. Only then is the local expansion integrated, with \emph{both} normalizing factors retained. The collision extension in Section~\ref{spm:sec:collision} uses the same three parts rather than a formal continuation of the final formula.

\section{A positive envelope and all nonreal poles}\label{spm:sec:envelope}
Choose $U$ small enough that $\Re L>c>0$, $|L|<1$, $\Re v>c>0$, and $|v|\leq V$ for a fixed real $V$. Fix a real radius $R$ satisfying
\begin{equation}\label{spm:eq:R}
 \sup_U|L|<R<\inf_{U,\,m\ne0}|L+2\pi\ii m|;
\end{equation}
for example $R=2$ is possible after shrinking $U$.

Write $f_n(k,v)=\sum_{j=1}^n b_{nj}(v)k^j$ for $n>0$. These coefficients need not be nonnegative: already
\[
 f_2(k,v)=\frac{v^2+1}{2}k^2+\frac{v^2-1}{2}k.
\]
In particular positivity would fail for real $0<v<1$. The following replacement is valid coefficient by coefficient.

\begin{lemma}[Positive majorant]\label{spm:lem:envelope}
Define $g_n(k,V)$ by
\begin{equation}\label{spm:eq:envelope}
 \sum_{n\geq0}g_n(k,V)z^n
 =(1-Vz)^{-k}(1-z^2)^{-(k^2+k)/2}.
\end{equation}
All coefficients $g_{nj}=[k^j]g_n(k,V)$ are nonnegative, and
\[
 |b_{nj}(v)|\leq g_{nj},\qquad
 |f_n(k,v)|\leq g_n(|k|,V).
\]
\end{lemma}
\begin{proof}
The logarithm of \eqref{spm:eq:fn} is
\[
 -k\log(1-vz)+\frac{k}{2}\log(1-z^2)
 -\frac{k^2}{2}\log(1-z^2).
\]
Take absolute values coefficientwise in $k,z$. Its coefficients are bounded by those of the logarithm of \eqref{spm:eq:envelope}; exponentiation preserves the coefficientwise bound. The majorant has nonnegative coefficients, so $g_n(t,V)$ is also increasing for $t\geq0$.
\end{proof}

For every integer $j\geq1$ and $\Re L>0$, the absolutely convergent partial-fraction identity is
\begin{equation}\label{spm:eq:poles}
 \sum_{k\geq0}e^{-Lk}k^j
 =j!\sum_{m\in\mathbb Z}(L+2\pi\ii m)^{-j-1}.
\end{equation}
One obtains it by differentiating the partial-fraction expansion of $(e^L-1)^{-1}$; one derivative removes the possible constant term and makes the series absolutely convergent. Separating the $m=0$ contribution therefore gives the \emph{finite polynomial} identity
\begin{align}
 \sum_{k\geq0}e^{-Lk}f_n(k,v)&=M_n(L,v)+E_n(L,v),\label{spm:eq:mainpluserror}\\
 M_n(L,v)&=L^{-1}\int_0^\infty e^{-t}f_n(t/L,v)\,\dd t.
 \label{spm:eq:gammaexact}
\end{align}
No integral of an untruncated infinite generating function has been interchanged here.

The nonreal terms admit a degree-independent bound. Indeed, if $d_m=|L+2\pi\ii m|$, then uniformly for $j\geq1$,
\[
 \sum_{m\ne0}d_m^{-j-1}
 \leq R^{-j-1}\sum_{m\ne0}(R/d_m)^2
 \leq C R^{-j-1}.
\]
Thus Lemma~\ref{spm:lem:envelope} yields
\begin{equation}\label{spm:eq:allpolebound}
 |E_n|\leq C\sum_{j=1}^n g_{nj}j!R^{-j-1}
 =C\int_0^\infty e^{-Rk}g_n(k,V)\,\dd k.
\end{equation}
This bound includes degrees $1$ and $2$, rather than incorrectly attaching an $n$th power of a pole ratio to each low-degree monomial.

\begin{lemma}[Global gamma and Cauchy bound]\label{spm:lem:gammabound}
For fixed positive $R,V$,
\begin{equation}\label{spm:eq:majorantintegral}
 \int_0^\infty e^{-Rk}g_n(k,V)\,\dd k
 \leq C\sqrt n\,I_n(V)R^{-n-1}.
\end{equation}
The analogous estimate is uniform when the nonzero scale parameter ranges over a fixed compact complex set and its absolute value is used in the envelope.
\end{lemma}
\begin{proof}
Put $k=t/|\lambda|$, where $\lambda$ stays in a fixed compact subset of $\C\setminus\{0\}$. For $t\geq\eps n$, take $|x|=\sqrt n$ and $z=x/k$. For sufficiently large $n$, all logarithmic series converge. After removing $Vx+x^2/2$ from the logarithm of \eqref{spm:eq:envelope}, the remaining terms are bounded by a constant depending on $\eps$ and the compact set. The largest types are $x^2/k$ and $x^4/k^2$, both $O(1)$. Cauchy's coefficient estimate gives
\begin{equation}\label{spm:eq:largekbound}
 g_n(t/|\lambda|,V)
 \leq C_\eps(t/|\lambda|)^n n^{-n/2}e^{n/2+V\sqrt n}.
\end{equation}
Integration against $e^{-t}$ uses $\int_0^\infty e^{-t}t^n\,\dd t=n!$.

For $0\leq t<\eps n$, use coefficient monotonicity before making any substitution involving $1/k$:
\[
 \int_0^{\eps n}e^{-t}g_n(t/|\lambda|,V)\,\dd t
 \leq \eps n\,g_n(\eps n/|\lambda|,V).
\]
Apply \eqref{spm:eq:largekbound} at the endpoint. Relative to
$n!|\lambda|^{-n}n^{-n/2}e^{n/2+V\sqrt n}$, the ratio is bounded by a polynomial in $n$ times $(e\eps)^n$. Taking $\eps<e^{-2}$ makes it exponentially small. This covers $k=0$ and arbitrarily small positive $k$.

The involution saddle established independently in Lemma~\ref{spm:lem:involution} below gives
\[
 I_n(V)\asymp n!n^{-n/2}e^{n/2+V\sqrt n}/\sqrt n.
\]
Comparing with the preceding bound and setting $\lambda=R$ proves \eqref{spm:eq:majorantintegral}. The factor $\sqrt n$ is harmless here and need not be optimized away.
\end{proof}

The same saddle gives, uniformly near $v=1$,
\[
 |I_n(v)|\geq c\,n!n^{-n/2}e^{n/2+\Re(v)\sqrt n}/\sqrt n.
\]
Consequently \eqref{spm:eq:allpolebound} implies
\begin{equation}\label{spm:eq:poleerror}
 |E_n|\leq C e^{-c_1 n}|I_n(v)L^{-n-1}|.
\end{equation}
Indeed its relative upper bound is a polynomial times
$e^{C\sqrt n}(|L|/R)^{n+1}$, exponentially small by \eqref{spm:eq:R}. This is the required simultaneous control of every polynomial degree.

\section{The two saddle variables and their tails}\label{sec:saddles}
\subsection{The complex involution contour}
Let
\begin{equation}\label{eq:r}
 r=r_n(v)=\sqrt{n+v^2/4}-v/2,\qquad r^2+vr=n,
\end{equation}
where $r/\sqrt n\to1$. Even when $r$ is complex, $x=re^{\ii\theta}$ parametrizes a permissible circle around the origin. Relative to $\theta=0$, the coefficient exponent is
\begin{equation}\label{eq:Phi}
 \Phi_n(\theta)=vr(e^{\ii\theta}-1)
 +\frac{r^2}{2}(e^{2\ii\theta}-1)-\ii n\theta.
\end{equation}
Its first derivative at zero vanishes exactly.

\begin{lemma}[Uniform angular localization]\label{lem:involution}
For a sufficiently small fixed complex neighborhood of $v=1$ and $0<\delta<1/6$,
\begin{equation}\label{eq:angulartail}
 \int_{|\theta|>n^{-1/2+\delta}}|e^{\Phi_n(\theta)}|\,\dd\theta
 =O\bigl(n^{-1/2}e^{-c n^{2\delta}}\bigr).
\end{equation}
The full absolute integral is $O(n^{-1/2})$. The signed integral is asymptotic to $\sqrt{\pi/n}$, uniformly, and hence
\begin{equation}\label{eq:involutioncomparability}
 |I_n(v)|\asymp
 \frac{n!}{\sqrt n}\left|r^{-n}e^{vr+r^2/2}\right|.
\end{equation}
In particular $I_n(v)$ is nonzero there for all sufficiently large $n$.
\end{lemma}
\begin{proof}
Using $r^2=n-vr$, rewrite the exponent as
\[
 \Phi_n(\theta)=\frac n2(e^{2\ii\theta}-1-2\ii\theta)
 +vr\left(e^{\ii\theta}-1-\frac{e^{2\ii\theta}-1}{2}\right).
\]
Its leading real part is $-n\sin^2\theta$. Near zero the expression multiplying $vr$ is $O(\theta^2)$, and $vr=O(\sqrt n)$, so $\Re\Phi_n\leq-cn\theta^2$ on a fixed small neighborhood. No spurious linear error remains.

Near $\pi$ or $-\pi$, put $\theta=\pi+\phi$ modulo $2\pi$. The constant term is $-2vr$, with $\Re(vr)\geq c\sqrt n$. The remaining terms are bounded by $C\sqrt n|\phi|-c'n\phi^2$. Completing the square shows that this saddle has height at most $-c\sqrt n+O(1)$ and Gaussian width $O(n^{-1/2})$. For complex $v$ its maximum may be displaced by $O(n^{-1/2})$; this changes the height only by $O(1)$. Away from both angular neighborhoods the bound is $-cn$. These three bounds prove \eqref{eq:angulartail} and the absolute-integral estimate.

On $\theta=w/\sqrt n$, Taylor expansion gives $\Phi_n=-w^2+o(1)$ uniformly on each bounded interval. The preceding tails give an integrable majorant, so the signed integral is $\sqrt{\pi/n}(1+o(1))$. The Cauchy formula proves \eqref{eq:involutioncomparability}. Expanding $r=\sqrt n-v/2+O(n^{-1/2})$ gives the size estimates used in Section~\ref{sec:envelope}, independently of its conclusions.
\end{proof}

\subsection{The gamma variable over its whole range}
For sufficiently large $k$ relative to $\sqrt n$, factor the scaled generating function as
\begin{equation}\label{eq:factor}
 (1-vx/k)^{-k}(1-x^2/k^2)^{-(k^2-k)/2}
 =e^{vx+x^2/2}e^{B(x,k,v)},
\end{equation}
where the convergent logarithmic correction is
\begin{align}\label{eq:B}
 B(x,k,v)={}&\sum_{m\geq2}\frac{v^m x^m}{m k^{m-1}}
 +\sum_{m\geq2}\frac{x^{2m}}{2m k^{2m-2}}
 -\sum_{m\geq1}\frac{x^{2m}}{2m k^{2m-1}}.
\end{align}
For $t\geq\eps n$, $k=t/L$, and $|x|=|r_n(v)|$, this correction is uniformly bounded in modulus. On that range \eqref{eq:gammaexact} becomes
\begin{equation}\label{eq:mainintegral}
 L^{-n-1}\int e^{-t}t^n
 [x^n]e^{vx+x^2/2+B(x,t/L,v)}\,\dd t.
\end{equation}
This is first an identity for a fixed coefficient polynomial, not an exchange with a divergent full generating-function integral.

The absolute angular bound and \eqref{eq:involutioncomparability} show that the coefficient inside the integral has absolute value at most $C|I_n(v)|/n!$. Thus, relative to $I_n(v)L^{-n-1}$, the mass outside
\begin{equation}\label{eq:gammawindow}
 |t-n|\leq n^{1/2+\delta}
\end{equation}
is $O(e^{-c n^{2\delta}})$ on $t\geq\eps n$. The gamma estimate follows directly from
\[
 \frac{t^ne^{-t}}{n^ne^{-n}}
 =\exp\left[-n\left(\frac tn-1-\log\frac tn\right)\right].
\]
The rate function $x-1-\log x$ is at least a positive constant times $(x-1)^2$ on a fixed interval about $1$ and is bounded away from zero off that interval, with sufficient growth to integrate the tails.

On $0\leq t<\eps n$, use Lemma~\ref{lem:envelope}, monotonicity, and the endpoint estimate from Lemma~\ref{lem:gammabound}. Dividing by \eqref{eq:involutioncomparability} gives
\[
 O\bigl((e\eps)^n e^{C\sqrt n}\operatorname{poly}(n)\bigr),
\]
which is exponentially small for fixed sufficiently small $\eps$. This separate treatment is essential: the substitution $x/k$ would be invalid at $k=0$.

Together these arguments localize the product gamma/angular integral with absolute estimates for complex markers. No positivity of a complex integrand is used.

\subsection{Real auxiliary and diagonal concentration}
The preceding bounds also explain the two scales in elementary probabilistic terms. For positive real $u,v$ in a fixed compact neighborhood, give the auxiliary integer $K$ probability proportional to $e^{-Lk}f_n(k,v)$. Uniformly for $cn\leq k\leq Cn$ and real $w$ near $v$, the same angular saddle yields
\begin{equation}\label{eq:centralcomparison}
 \frac{f_n(k,w)}{k^n I_n(w)/n!}
 =e^{B(r_n(w),k,w)}(1+O(n^{-1/2})).
\end{equation}
The multiplier is positive and bounded away from zero. Its change on the angular scale is $O(n^{-1/2})(1+|w_{\rm ang}|)$, which proves the error after Gaussian integration. For $k\geq\eps n$ the corresponding upper bound holds on the entire range. A central interval supplies a matching lower normalization. Comparison of the unimodal kernel $e^{-Lk}k^n$ with adjacent integrals then gives
\begin{equation}\label{eq:Ktail}
 \PP\{|K-n/L|>n^{1/2+\delta}\}
 \leq C e^{-c n^{2\delta}},\qquad 0<\delta<1/6.
\end{equation}
The small-$k$ range is covered by the majorant with $\eps$ chosen relative to the compact $L$ range. The auxiliary $K$ is not the matrix dimension $D$.

For the reference involution distribution, let $\Lambda$ have probabilities proportional to
\[
 \frac{v^\ell}{\ell!2^j j!},\qquad \ell+2j=n.
\]
The one-dimensional saddle, expanded with uniform derivatives, gives for real small $s$
\begin{equation}\label{eq:reference-mgf}
 \log\frac{I_n(ve^s)}{I_n(v)}
 =v(e^s-1)\sqrt n-\frac{v^2(e^{2s}-1)}4+O(n^{-1/2}).
\end{equation}
Chernoff's inequality with $s$ proportional to $n^{-1/4+\delta}$ therefore yields
\begin{equation}\label{eq:fixedpointtail}
 \PP\{|\Lambda-v\sqrt n|>n^{1/4+\delta}\}
 \leq C e^{-c n^{2\delta}},\qquad 0<\delta<1/4.
\end{equation}
This is not merely a reference-distribution bound. By \eqref{eq:centralcomparison}, the tilted diagonal generating-function ratio for actual matrices at each $cn\leq k\leq Cn$ is at most a constant times $I_n(ve^s)/I_n(v)$. Hence their diagonal sum obeys the same bound, with changed constants. This explicitly controls the rising-factorial reweighting that could otherwise invalidate a fixed-point moment argument. The complex-marker proof above does not depend on these probabilistic comparisons.

\section{Finite coefficient construction and rigorous remainder}\label{spm:sec:allorders}
Set $h=n^{-1/2}$ and make the central substitutions
\begin{equation}\label{spm:eq:scaling}
 t=h^{-2}(1+hy),\quad \theta=hw,\quad
 \rho(h)=\sqrt{1+v^2h^2/4}-vh/2,
\end{equation}
so that
\[
 x=h^{-1}\rho(h)e^{\ii hw},\qquad
 k=h^{-2}(1+hy)/L.
\]
Let $\mathcal E$ be the product Gaussian functional with independent variables $Y,W$ satisfying
\[
 Y\sim N(0,1),\qquad
 \PP(W\in\dd w)=\pi^{-1/2}e^{-w^2}\,\dd w.
\]
It is enough to define it on polynomials:
\[
 \mathcal E(Y^{2a})=(2a-1)!!,\qquad
 \mathcal E(W^{2b})=\frac{(2b-1)!!}{2^b},
\]
with zeroth moments $1$, odd moments $0$, and products factorized.

Define formal series with removable singularities at $h=0$ by
\begin{align}
 G(h,y)&=h^{-2}\bigl(\log(1+hy)-hy\bigr)+y^2/2,\label{spm:eq:G}\\
 H(h,w)&=v h^{-1}\rho(h)(e^{\ii hw}-1)
 +\frac{\rho(h)^2}{2h^2}(e^{2\ii hw}-1)-\frac{\ii w}{h}+w^2,\label{spm:eq:H}\\
 \mathcal B(h,y,w)&=B\bigl(h^{-1}\rho(h)e^{\ii hw},
                       h^{-2}(1+hy)/L,v\bigr).\label{spm:eq:scaledB}
\end{align}
Here $G(0,y)=H(0,w)=0$ and $\mathcal B(0,y,w)=S(L,v)$.

\begin{proposition}[Finite Gaussian algorithm]\label{spm:prop:algorithm}
The coefficients in Theorem~\ref{spm:thm:main} are defined by
\begin{equation}\label{spm:eq:algorithm}
 \sum_{j\geq0}C_j(L,v)h^j
 =e^{-S(L,v)}
 \frac{\mathcal E\exp(G+H+\mathcal B)}
      {(\mathcal E\exp G)(\mathcal E\exp H)}.
\end{equation}
This is an identity of formal power series: at any fixed order only finitely many Gaussian moments are taken. Both denominator factors are required.
\end{proposition}
\begin{proof}
On the central window, the gamma exponent is $-y^2/2+G$ and the angular exponent is $-w^2+H$. Dividing by the exact gamma normalization and the exact involution coefficient integral gives respectively the two factors in the denominator. In particular the gamma density is centered at $n$, not its mean $n+1$; the series $G$ and its normalizer retain that shift.

The least powers of $h$ in the three summand families of \eqref{spm:eq:B} are respectively $m-2$, $2m-4$, and $2m-2$. To calculate through $h^J$ one therefore needs only
\[
 m\leq J+2,\qquad m\leq\lfloor J/2\rfloor+2,
 \qquad m\leq\lfloor J/2\rfloor+1
\]
in those families. Each surviving coefficient is a polynomial in $y,w,L,v$ with rational coefficients apart from powers of $\ii$. Gaussian evaluation kills odd powers of $w$; the surviving powers of $\ii$ are real. Denominator division, with both constant terms equal to $1$, proves the rational-polynomial assertion.
\end{proof}

For transparency, the first calculation can be displayed without a symbolic system. Put $a=(v^2-1)/2$. The first nonconstant terms are
\begin{align*}
 G_1&=y^3/3,\qquad H_1=vw^2/2-2\ii w^3/3,\\
 \mathcal B_1&=aL(-v+2\ii w-y)
       +\frac{L^2}{4}(-2v+4\ii w-2y)+\frac{v^3L^2}{3}.
\end{align*}
The normalizers cancel the expectations of $G_1,H_1$, leaving
$C_1=\mathcal E\mathcal B_1$, exactly \eqref{spm:eq:C1}. At second order,
\begin{equation}\label{spm:eq:secondalgorithm}
 C_2=\mathcal E\left(\mathcal B_2+(G_1+H_1)\mathcal B_1
                         +\frac{\mathcal B_1^2}{2}\right)
       -\mathcal E(G_1+H_1)C_1.
\end{equation}
Substitution in this finite polynomial identity gives \eqref{spm:eq:D2}--\eqref{spm:eq:C2}. The reproduction package checks these identities exactly and separately checks the collision first correction.

\subsection{Why the formal algorithm is asymptotic}
A formal Gaussian calculation alone would not prove the theorem. Choose a fixed $0<\alpha<1/3$ and restrict $|y|,|w|\leq h^{-\alpha}$. The previous section bounds the excluded tails by $O(e^{-c h^{-2\alpha}})$. On this window, $1+hy$ is separated from zero, $\rho$ is analytic, and $|x/k|=O(h)$. The logarithmic series \eqref{spm:eq:B} and any fixed number of its $h$ derivatives converge uniformly.

Taylor's theorem through order $J$ for the local exponential has an integrable remainder bounded by
\begin{equation}\label{spm:eq:taylorbound}
 C_Jh^{J+1}(1+|y|+|w|)^{D_J}e^{\eta(y^2+w^2)},
 \qquad \eta<1/4.
\end{equation}
Here is a direct justification valid also at the intermediate points in Taylor's integral remainder. For every $0\leq s\leq h$, the identity
$\rho(s)^2+vs\rho(s)=1$ removes the apparent singularity in $H$. Uniformly on the window,
\[
 |G(s,y)|\leq Cs|y|^3,\qquad
 |H(s,w)|\leq Cs(|w|^2+|w|^3),\qquad
 |\mathcal B(s,y,w)|\leq C.
\]
Because $h(|y|+|w|)=o(1)$, the residual real parts are at most
$\eta(y^2+w^2)+C$ after taking $h$ sufficiently small. Every fixed derivative is bounded by a polynomial in $|y|+|w|$: the only denominator factors are powers of $1+sy$ and analytic factors separated from zero, and the least-degree structure of \eqref{spm:eq:B} bounds the number of relevant terms. Differentiating the exponential a fixed number of times multiplies it only by such polynomials. This proves \eqref{spm:eq:taylorbound}.

Multiplication by the Gaussian density $e^{-y^2/2-w^2}$ makes \eqref{spm:eq:taylorbound} integrable over the full plane. Its integral is $O_J(h^{J+1})$. Extending each finite polynomial coefficient integral from the window to the full plane loses only an exponentially small tail. The same argument applies to both denominator integrals, whose constant terms are nonzero. This proves \eqref{spm:eq:main} after the exponentially small pole error \eqref{spm:eq:poleerror} is restored.

The notation $\mathcal E e^{G+H+\mathcal B}$ in \eqref{spm:eq:algorithm} is formal. The proof extends only finite polynomial truncations to the full real plane; it never integrates $\log(1+hy)$ through $y=-1/h$.

Finally the normalized quantity in \eqref{spm:eq:main} is holomorphic and nonzero on a slightly smaller neighborhood for large $n$, by Lemma~\ref{spm:lem:involution} and its uniform limit $1$. Cauchy's formula, applied on a larger compact polydisc, bounds all fixed derivatives of the remainder. Its analytic logarithm consequently has a differentiable expansion to every fixed order. This completes the proof of Theorem~\ref{spm:thm:main}.

\begin{remark}[{[write]} The third coefficient at $u=v=1$, and the data, 6~October 2026]\label{spm:rem:third}
The write evaluated the finite prescription~\eqref{spm:eq:algorithm} at
$u=v=1$ through $h^3$ in SymPy~1.14.0, with $L$ symbolic (an implementation
of its own, not shipped; it expands $G$, $H$ and $\mathcal B$ from
\eqref{spm:eq:G}--\eqref{spm:eq:scaledB}, keeping in~\eqref{spm:eq:B}
exactly the terms listed in the proof of
Proposition~\ref{spm:prop:algorithm}, takes the Gaussian moments of
$\mathcal E$, and divides by both normalizers). It returns $S(L,1)=L^2/4$,
$C_1(L,1)=-L^2/6$ and $C_2(L,1)=-L^2/4+L^4/18$, as in
\eqref{spm:eq:C1}--\eqref{spm:eq:C2}, and
\[
 C_3(L,1)=-\frac{L^2(50L^4+324L^2-1755)}{6480},\qquad
 D_3(L,1)=C_3-C_1C_2+\tfrac13C_1^3=-\frac{L^2(22L^2-65)}{240}.
\]
Theorem~\ref{spm:thm:main} with $J=3$ therefore extends~\eqref{spm:eq:unmarked}
by the term $C_3n^{-3/2}$, with remainder $O(n^{-2})$; at $L=\log2$,
$C_1=-0.08007550\ldots$, $C_2=-0.10728908\ldots$ and
$C_3=0.11772518\ldots$. This value rests on the write's single
implementation (Question~\ref{spm:q:third}).

\emph{The data.} With $\varrho_n=A_n\big/\bigl(e^{L^2/4}I_n(1)/(2L^{n+1})\bigr)$
and $h=n^{-1/2}$, the exact counts of Remark~\ref{spm:rem:oeis} give
(50-digit \texttt{mpmath}; the delivered noncertified receipt
\file{data/diagnostic_receipt.json} records the quantities of the second to fifth columns at
$n=32$, $80$, $160$, $320$, $640$, and agrees with the row $n=32$):
\begin{center}\small
\begin{tabular}{@{}r*{5}{c}@{}}
\toprule
$n$ & $\varrho_n$ & $(\varrho_n-1)/h$ & $(\varrho_n-1-C_1h)/h^2$ & $(\varrho_n-1-C_1h-C_2h^2)/h^3$ & $(\varrho_n-\sum_{j\le3}C_jh^j)/h^4$\\
\midrule
16 & 0.97438 & $-0.1025$ & $-0.0896$ & 0.0709 & $-0.1872$\\
32 & 0.98295 & $-0.0965$ & $-0.0927$ & 0.0826 & $-0.1986$\\
50 & 0.98678 & $-0.0935$ & $-0.0947$ & 0.0889 & $-0.2041$\\
100 & 0.99102 & $-0.0898$ & $-0.0976$ & 0.0966 & $-0.2112$\\
200 & 0.99384 & $-0.0872$ & $-0.1000$ & 0.1024 & $-0.2168$\\
300 & 0.99504 & $-0.0859$ & $-0.1012$ & 0.1051 & $-0.2194$\\
400 & 0.99574 & $-0.0852$ & $-0.1020$ & 0.1067 & $-0.2210$\\
500 & 0.99621 & $-0.0847$ & $-0.1025$ & 0.1078 & $-0.2222$\\
\bottomrule
\end{tabular}
\end{center}
(rounded to the digits shown). The third to fifth columns move toward
$C_1$, $C_2$ and $C_3$ as $n$ grows, and the last stays between $-0.18$
and $-0.23$ for $16\le n\le500$. These are observations on a finite range,
consistent with the expansion; they certify no coefficient and no rate.
\end{remark}
\begin{writenote}[Independent check of the write, 6~October 2026]
An adversarial check made by the intake after the write (commit
\file{c569b0b98}) re-derived the statements marked [write] with its own
code, after fetching again A138178, its b-file and A135588 (the revisions
quoted) and Cameron--Prellberg--Stark's v2 (4~April 2006).
\emph{Counts.} By a route different from the write's, the zero-line
inclusion--exclusion~\eqref{spm:eq:deletion} with $f_n(j)$ from the
first-order recurrence
$(n+1)a_{n+1}=ja_n+(n-1+j+j(j-1))a_{n-1}$ of
$(1-z)^{-j}(1-z^2)^{-j(j-1)/2}$, it computed $A_0,\ldots,A_{2000}$ exactly:
all 501 b-file terms and the 25 data terms agree; binomial cell factors
give all 26 terms of A135588. \emph{Remark~\ref{spm:rem:third}.} Every entry
of the table was reproduced to the digits shown, and the receipt's row
$n=32$ agrees. Richardson extrapolation in $h$ of
$(\varrho_n-1-C_1h-C_2h^2)/h^3$ from the exact counts gives
$0.1177251847$ (nodes $n=400,600,\ldots,2000$; the same ten digits from
three other node sets of seven to thirteen points), against $C_3(\log2,1)=0.1177251846518\ldots$ from the
write's formula; the same extrapolation of marked counts $A_n(u,1)$ to
$n=1500$ gives $0.154285415$ at $u=\tfrac45$ and $0.087284579$ at
$u=\tfrac54$, against $0.1542854152\ldots$ and $0.0872845791\ldots$ from the
formula at $L=\log\tfrac94$ and $L=\log\tfrac95$. This is numerical
confirmation at three values of $L$, independent of the Gaussian
prescription; it is not a derivation (Question~\ref{spm:q:third}).
$D_3=C_3-C_1C_2+\tfrac13C_1^3$ and $C_1(L,1)$, $C_2(L,1)$ were rechecked in
SymPy. \emph{Remark~\ref{spm:rem:oeis}.} The entries' revisions, names,
generating functions (the first is~\eqref{spm:eq:deletion} at $u=v=1$,
Jovovic's is~\eqref{spm:eq:transform}), comments and the absence of an
asymptotic are as quoted; Proposition~4.2 (p.~14, with $C_s\approx0.44341$)
and the Poisson$((\log2)^2/2)$ limit with~(8) (p.~9) are in the preprint as
described. \emph{Section~\ref{spm:sec:model}, note.} \file{q2:thm:fubini},
Part~IV of \file{a261781-matrix-compositions} (Sections~34--44, the polydisc
$|\omega-1|\le2.1$, $|u-1|\le0.02$, $P\sim\mathrm{Pois}(Lr)$ with
$r\log2=1.1046161\ldots$) and the Lean names are as stated.
\emph{Section~\ref{spm:sec:saddles}, note.} The identity in law was
re-derived ($\sum_{k\ge d}q^k\binom kd=u^d/(1-q)$, and $K$ has the law
proportional to $e^{-Lk}f_n(k,v)=q^kf_n(k,v)$ there).
\emph{Remark~\ref{spm:rem:transseries}.} (1)~holds; (2)~the factorial core
with $\kappa=\tfrac12$, $d=-\tfrac12-\log L$ gives $v_{\rm vol}=w$,
$X_{\rm vol}=N$, slope $Q/2$, and the core equation alone has the
\file{plt:thm:lw-template} data stated, the extra terms being of relative
order $e^{-Z/2}/Z$; (3)~the identity for $F_0$ is exact, the triangular
recurrence gives $e_1,e_2,e_3=\alpha,\beta,\gamma$, and
$e_1,\dots,e_5$ are polynomials in $Q^{-1}$ without constant term, of
degrees $1,3,5,7,9$, with $d$ first in $e_4$ (SymPy); (4)~is accurate.
\emph{Provenance.} The archive facts (604{,}411 bytes, the SHA-256 quoted,
30 files, 29 manifest entries, 952 delivered lines, 22 pages) and the counts
of 110 labels and 93 references were confirmed. No defect was found.
\end{writenote}

\section{Dimension and trace fluctuations}\label{spm:sec:fluctuations}
Choose a matrix uniformly from $\mathcal A_n$, and write $D_n,T_n$ for its dimension and diagonal sum. Their joint moment generating function is
\[
 \E e^{aD_n+bT_n}=\frac{A_n(e^a,e^b)}{A_n(1,1)}.
\]
Theorem~\ref{spm:thm:main} applies in a fixed complex neighborhood of $(a,b)=(0,0)$, including derivative remainders. It therefore supplies all fixed cumulants by differentiation.

\subsection{The involution factor}
The one-dimensional saddle of Lemma~\ref{spm:lem:involution}, with Stirling's formula for $n!$, gives uniformly near $v=1$
\begin{align}\label{spm:eq:logI}
 \log I_n(v)={}&\frac n2(\log n-1)+v\sqrt n-\frac{v^2}{4}
              -\frac12\log2\notag\\
 &+\left(\frac{v^3}{24}+\frac v4\right)n^{-1/2}
   -\left(\frac{v^2}{16}+\frac1{12}\right)n^{-1}
   +O(n^{-3/2}).
\end{align}
For completeness, expand $\rho(h)$ from \eqref{spm:eq:scaling} in
$vr+r^2/2-n\log r$, expand the angular normalizer using $H$ in \eqref{spm:eq:H}, and add
\[
 \log n!=(n+\tfrac12)\log n-n+\tfrac12\log(2\pi)
             +\frac1{12n}+O(n^{-3}).
\]
At each order this uses finitely many even Gaussian moments. The same Taylor bound as \eqref{spm:eq:taylorbound}, with no gamma variable, proves the remainder. Thus \eqref{spm:eq:logI} is also an all-fixed-order construction. At $v=1$ its logarithmic coefficient of $n^{-1}$ is $-7/48$.

\begin{proposition}[First joint cumulants]\label{spm:prop:cumulants}
With $L=\log2$,
\begin{align}
 \E D_n={}&\frac{n}{2L}+\frac1{2L}-\frac12-\frac L4+O(n^{-1/2}),\label{spm:eq:ED}\\
 \Var D_n={}&\frac{n(1-L)}{4L^2}
 +\frac{1-L}{4L^2}+\frac{L-1}{8}+O(n^{-1/2}),\label{spm:eq:VD}\\
 \E T_n={}&\sqrt n+L-\frac12
 +\frac{3/8-L+L^2/2}{\sqrt n}+O(n^{-1}),\label{spm:eq:ET}\\
 \Var T_n={}&\sqrt n+2L-1
 +\frac{5/8-4L+5L^2/2}{\sqrt n}+O(n^{-1}),\label{spm:eq:VT}\\
 \Cov(D_n,T_n)={}&-\frac12+\frac{1-L}{2\sqrt n}+O(n^{-1}).\label{spm:eq:cov}
\end{align}
\end{proposition}
\begin{proof}
Take the logarithm of \eqref{spm:eq:main}, using
$C_1n^{-1/2}+D_2n^{-1}$ for its first normalized logarithmic terms, and insert \eqref{spm:eq:logI}. Differentiate with respect to $a$ and $b$. The required identities are
\[
 \left.\frac{\dd L(e^a)}{\dd a}\right|_{a=0}=-\frac12,
 \qquad
 \left.\frac{\dd^2 L(e^a)}{\dd a^2}\right|_{a=0}=\frac14,
 \qquad \partial_b=v\partial_v.
\]
Both $-(n+1)\log L(e^a)$ and $-\log(1+e^a)$ must be retained. The uniform holomorphic remainder justifies the differentiated error bounds.
\end{proof}

\begin{theorem}[Two independent Gaussian scales]\label{spm:thm:gaussian}
Let $\sigma_D^2=(1-L)/(4L^2)>0$. Then
\begin{equation}\label{spm:eq:gaussian}
 \left(\frac{D_n-n/(2L)}{\sigma_D\sqrt n},
       \frac{T_n-\sqrt n}{n^{1/4}}\right)
 \ \Longrightarrow\ (Z_1,Z_2),
\end{equation}
where $Z_1,Z_2$ are independent standard normal variables.
\end{theorem}
\begin{proof}
Substitute $a=s/(\sigma_D\sqrt n)$ and $b=t/n^{1/4}$ in the logarithm of the marked ratio. After centering, its limit is $(s^2+t^2)/2$ on a neighborhood of the origin. The leading $a$ term comes from $-(n+1)\log L(e^a)$, the leading $b$ term from $e^b\sqrt n$, and all mixed terms vanish on these scales. Convergence of the joint moment generating function proves the result.
\end{proof}

The nonzero limiting \emph{unscaled} covariance in \eqref{spm:eq:cov} is compatible with independent scaled limits. There is also an exact lattice constraint:
\begin{equation}\label{spm:eq:parity}
 T_n\equiv n\pmod2.
\end{equation}
No span-one local limit or total-variation approximation of $T_n$ by an unrestricted Poisson law is asserted. The valid local analytic residue is instead
\begin{equation}\label{spm:eq:trace-residue}
 \E[v^{T_n}]e^{-(v-1)\sqrt n}
 \longrightarrow \exp\left(\frac{(2L-1)(v^2-1)}4\right)
\end{equation}
uniformly near $v=1$. This follows by combining the $v$-dependent terms in \eqref{spm:eq:logI} with $S(L,v)-S(L,1)$.

\section{Collision markers on arbitrary fixed compact sets}\label{sec:collision}
For a symmetric packed matrix, let $R_n$ count diagonal cells with entry at least $2$, and let $S_n$ count upper-triangular off-diagonal cells with entry at least $2$. Each symmetric off-diagonal pair is counted once. These statistics count repeated \emph{cells}, not the number of extra units of weight. Define
\[
 A_n(u,v;w,\zeta)=\sum_{M\in\mathcal A_n}
 u^{D(M)}v^{T(M)}w^{R(M)}\zeta^{S(M)}.
\]
A diagonal cell and an upper-triangular off-diagonal cell contribute respectively
\[
 1+vz+\frac{w(vz)^2}{1-vz}
   =\frac{1+(w-1)v^2z^2}{1-vz},\qquad
 1+z^2+\frac{\zeta z^4}{1-z^2}
   =\frac{1+(\zeta-1)z^4}{1-z^2}.
\]
Consequently, before zero-line deletion, the coefficient polynomial is
\begin{align}\label{eq:collision-fn}
 f_n(k,v;w,\zeta)=[z^n]{}&(1-vz)^{-k}(1-z^2)^{-(k^2-k)/2}\notag\\
 &\mathrel{}\times[1+(w-1)v^2z^2]^k
 [1+(\zeta-1)z^4]^{(k^2-k)/2}.
\end{align}
It still has degree at most $n$ in $k$, and zero constant term for $n>0$, since each logarithmic atom $k^az^b$ has $a\leq b$. Proposition~\ref{prop:transform} applies unchanged.

\begin{theorem}[Compact collision-marker expansion]\label{thm:collision}
There is a sufficiently small complex neighborhood $U$ of $(u,v)=(1,1)$ such that, for every fixed compact set $K\subset\C^2$ of collision markers $(w,\zeta)$ and every fixed integer $J\geq0$,
\begin{align}\label{eq:collision-theorem}
 A_n(u,v;w,\zeta)
 =\frac{I_n(v)e^{S_*(L,v,w,\zeta)}}{(1+u)L^{n+1}}
 \left(\sum_{j=0}^{J}C_j^*(L,v,w,\zeta)n^{-j/2}
       +O_{J,K}(n^{-(J+1)/2})\right)
\end{align}
uniformly on compact subsets of $U\times K$. Here $C_0^*=1$,
\begin{equation}\label{eq:Sstar}
 S_*=S(L,v)+(w-1)v^2L+(\zeta-1)L^2/2,
\end{equation}
and every $C_j^*$ is a rational-coefficient polynomial in $L,v,w,\zeta$. In particular $K$ may contain $(0,0)$.
\end{theorem}
\begin{proof}
We verify every required global and local modification of the proof of Theorem~\ref{thm:main}. The additional logarithmic correction on the scaled contour is
\begin{align}\label{eq:deltaB}
 \Delta B={}&k\log\left(1+(w-1)v^2x^2/k^2\right)\notag\\
 &+\frac{k^2-k}{2}\log\left(1+(\zeta-1)x^4/k^4\right).
\end{align}
For $t\geq\eps n$, $k=t/L$, and $|x|=|r_n(v)|$, the two arguments are $1+O_K(n^{-1})$ and $1+O_K(n^{-2})$. Their branches are uniquely specified by expansion at zero, and $\Delta B$ is uniformly bounded. On the central scale it tends to
$(w-1)v^2L+(\zeta-1)L^2/2$.

For the global envelope choose real $V\geq|v|$, $M\geq|w-1|$, and $N\geq|\zeta-1|$ uniformly on the fixed compact set. Replace \eqref{eq:envelope} by
\begin{align}\label{eq:collision-envelope}
 \sum_{n\geq0}g_n(k)z^n={}&(1-Vz)^{-k}(1-z^2)^{-(k^2+k)/2}\notag\\
 &\mathrel{}\times(1-MV^2z^2)^{-k}
 (1-Nz^4)^{-(k^2+k)/2}.
\end{align}
Expansion of the two new logarithms majorizes the absolute values of their counterparts in \eqref{eq:collision-fn}; replacing $k^2-k$ by $k^2+k$ also majorizes in the variable $k$. Exponentiation preserves the bound. All coefficients of this new envelope in $k,z$ are nonnegative.

For $k\geq\eps n$ and $|x|=\sqrt n$, the logarithm of \eqref{eq:collision-envelope} after removing $Vx+x^2/2$ is bounded by a compact-dependent constant. Its new largest terms are $MV^2x^2/k=O_K(1)$ and $Nx^4/(2k^2)=O_K(1)$. Hence the gamma/Cauchy bound remains valid. For small $k$, monotonicity gives the same endpoint estimate. Applying the exact partial-fraction formula to all positive powers of $k$ and using this envelope proves an exponentially small nonreal-pole error exactly as in \eqref{eq:allpolebound}--\eqref{eq:poleerror}. This includes every low degree.

On the involution contour, multiplication by $e^{\Delta B}$ changes the angular absolute bounds only by a fixed factor. The opposite saddle remains suppressed by $\Re v>0$. The local multiplier $e^{S_*}$ never vanishes, even at zero collision markers. Thus neither positivity of $w,\zeta$ nor an asymptotic continuation from their neighborhood of $1$ is needed.

Finally the scaled $\Delta B$ has a removable singularity at $h=0$, polynomial coefficients, and polynomially bounded fixed derivatives in the Gaussian variables. The argument for \eqref{eq:taylorbound} applies verbatim. Replace $\mathcal B$ by $\mathcal B+\Delta\mathcal B$ and $e^{-S}$ by $e^{-S_*}$ in \eqref{eq:algorithm}. This constructs every $C_j^*$, with both normalizers retained, and proves the stated remainder directly on each fixed compact set.
\end{proof}

A useful explicit coefficient is obtained by setting
\[
 a=\frac{v^2-1}{2}+(w-1)v^2,\qquad c=\frac{2\zeta-1}{4}.
\]
Then
\begin{equation}\label{eq:C1star}
 C_1^*=-v(aL+2cL^2)+\frac{v^3L^2}{3}.
\end{equation}
Indeed the leading correction atoms are $ax^2/k+cx^4/k^2$; the first-order Gaussian calculation in Section~\ref{sec:allorders} with these replacements gives \eqref{eq:C1star}. At $w=\zeta=1$ it recovers $C_1$.

\subsection{The binary specialization and its attribution}
Setting $w=\zeta=0$ forbids every entry of size at least $2$, so this specialization is exactly the symmetric binary model. Its amplitude and first correction are
\begin{align}
 S_*(L,v,0,0)&=-\frac{(v^2+1)L}{2}-\frac{L^2}{4},\label{eq:binaryS}\\
 C_1^*(L,v,0,0)&=\frac{v(v^2+1)L}{2}
                   +\left(\frac v2+\frac{v^3}{3}\right)L^2.\label{eq:binaryC1}
\end{align}
At $u=v=1$ this gives
\begin{equation}\label{eq:binaryrefined}
 B_n=\frac{e^{-L-L^2/4}I_n(1)}{2L^{n+1}}
 \left(1+\frac{L+5L^2/6}{\sqrt n}+O(n^{-1})\right).
\end{equation}
The leading term is precisely the established Cameron--Prellberg--Stark equivalent \eqref{eq:priorbinary}, recovered here as a consistency check. In particular
\begin{equation}\label{eq:binaryprob}
 \frac{B_n}{A_n}
 =e^{-L-L^2/2}\left(1+\frac{L+L^2}{\sqrt n}+O(n^{-1})\right).
\end{equation}
The probability of being binary has a nonzero limit even though the dimension diverges.

\section{Independent collisions and sharp total variation}\label{spm:sec:TV}
Throughout this section $L=\log2$ and
\[
 \lambda=L,\qquad \mu=L^2/2,\qquad \kappa=\lambda+2\mu=L+L^2.
\]

\begin{theorem}[Joint Gaussian and Poisson limits]\label{spm:thm:joint}
For every fixed compact set of $(w,\zeta)\in\C^2$,
\begin{equation}\label{spm:eq:collisionlimit}
 \E[w^{R_n}\zeta^{S_n}]
 \longrightarrow e^{\lambda(w-1)+\mu(\zeta-1)}
\end{equation}
uniformly. Moreover,
\begin{equation}\label{spm:eq:fourjoint}
 \left(\frac{D_n-n/(2L)}{\sigma_D\sqrt n},
       \frac{T_n-\sqrt n}{n^{1/4}},R_n,S_n\right)
 \Longrightarrow (Z_1,Z_2,P_1,P_2),
\end{equation}
where all four variables on the right are independent, $Z_1,Z_2$ are standard normal, $P_1\sim\Pois(\lambda)$, and $P_2\sim\Pois(\mu)$.
\end{theorem}
\begin{proof}
Normalize Theorem~\ref{spm:thm:collision} by its value at all markers equal to $1$. At $u=v=1$, its leading multiplier is the right side of \eqref{spm:eq:collisionlimit}. For the full joint limit substitute the scaled dimension and trace markers from Theorem~\ref{spm:thm:gaussian} and take $w=e^r,\zeta=e^q$ near $1$. Since $L(u)\to L$ and $v\to1$, the limiting log moment generating function is
\[
 (s^2+t^2)/2+\lambda(e^r-1)+\mu(e^q-1).
\]
It separates into four cumulant generating functions, proving the assertion. The finite-$n$ trace parity in \eqref{spm:eq:parity} is consistent with this weak limit.
\end{proof}

\subsection{A summable first correction}
Let $p_n(r,s)=\PP(R_n=r,S_n=s)$ and let
\[
 p_{r,s}=e^{-\lambda-\mu}\frac{\lambda^r\mu^s}{r!s!}
\]
be the independent Poisson mass function. From \eqref{spm:eq:C1star},
\[
 C_1^*(L,1,w,\zeta)-C_1(L,1)
 =\lambda(1-w)+2\mu(1-\zeta).
\]
Consequently the generating-function correction extracts to
\begin{equation}\label{spm:eq:masscorrection}
 p_n(r,s)=p_{r,s}\left[1+\frac{\kappa-r-2s}{\sqrt n}\right]+e_n(r,s),
 \qquad \sum_{r,s\geq0}|e_n(r,s)|=O(n^{-1}).
\end{equation}
The bracket is a signed correction; it is not claimed positive for every $r,s$.

The summability of the remainder is important. Fix any $\rho>1$. The compact-marker theorem applies on the entire closed polydisc $|w|,|\zeta|\leq\rho$. After scalar normalization, its additive generating-function remainder $E_n(w,\zeta)$ has supremum $O_\rho(n^{-1})$ there. Cauchy's estimate gives
\[
 |[w^r\zeta^s]E_n|\leq C_\rho n^{-1}\rho^{-r-s}.
\]
Summing the geometric bound proves \eqref{spm:eq:masscorrection}. Uniformity only near marker $1$ would not suffice for this argument; the larger compact domain was proved directly in Section~\ref{spm:sec:collision}.

\begin{theorem}[Sharp total-variation constant]\label{spm:thm:TV}
Use $\TV(P,Q)=\frac12\sum|P-Q|$. Then
\begin{equation}\label{spm:eq:sharpTV}
 \TV\bigl(\Law(R_n,S_n),\Pois(L)\otimes\Pois(L^2/2)\bigr)
 =\frac{e^{-L-L^2/2}L^2(2+L)}{\sqrt n}+O(n^{-1}).
\end{equation}
For $n>4$, set
\begin{equation}\label{spm:eq:adjustedmeans}
 \lambda_n=L(1-n^{-1/2}),\qquad
 \mu_n=\frac{L^2}{2}(1-2n^{-1/2}).
\end{equation}
These positive means give the genuine probability approximation
\begin{equation}\label{spm:eq:correctedTV}
 \TV\bigl(\Law(R_n,S_n),\Pois(\lambda_n)\otimes\Pois(\mu_n)\bigr)
 =O(n^{-1}).
\end{equation}
\end{theorem}
\begin{proof}
Under $p$, $\E_p(\kappa-R-2S)=0$. Thus \eqref{spm:eq:masscorrection} and the inequality $||x+y|-|x||\leq|y|$ give
\[
 \TV(\Law(R_n,S_n),p)
 =n^{-1/2}\E_p[(\kappa-R-2S)_+]+O(n^{-1}).
\]
Since $1<\kappa<2$, only $(R,S)=(0,0)$ and $(1,0)$ contribute to the positive part. Therefore
\begin{align*}
 \E_p[(\kappa-R-2S)_+]
 &=e^{-\lambda-\mu}\bigl[\kappa+\lambda(\kappa-1)\bigr]\\
 &=e^{-L-L^2/2}L^2(2+L),
\end{align*}
proving \eqref{spm:eq:sharpTV}, including its factor of $1/2$.

The pgf of the adjusted Poisson pair is
\[
 e^{\lambda(w-1)+\mu(\zeta-1)}
 \exp\left(n^{-1/2}[\lambda(1-w)+2\mu(1-\zeta)]\right).
\]
Its expansion has the same first correction, with uniform $O(n^{-1})$ remainder on the same radius-$\rho$ polydisc. The same Cauchy summation proves \eqref{spm:eq:correctedTV}.
\end{proof}

No optimality among all possible approximating distributions is claimed. Neither the $O(n^{-1})$ constant nor a numerical starting index has been made effective.

\subsection{Every fixed order in the joint mass law}
There is also a useful general consequence, requiring no new saddle analysis. Form the formal normalized coefficient series
\begin{equation}\label{spm:eq:Qpolynomials}
 \sum_{j\geq0}Q_j(w,\zeta)h^j
 =\frac{\sum_{j\geq0}C_j^*(L,1,w,\zeta)h^j}
        {\sum_{j\geq0}C_j(L,1)h^j}.
\end{equation}
Each $Q_j$ is a finite polynomial. If $Q_j=\sum_{a,b}q_{j,a,b}w^a\zeta^b$, define
\begin{equation}\label{spm:eq:Charlier}
 H_j(r,s)=\sum_{a,b}q_{j,a,b}
        \frac{(r)_a(s)_b}{\lambda^a\mu^b},
\end{equation}
where $(r)_a=r(r-1)\cdots(r-a+1)$ and $(r)_0=1$. For nonnegative integer $r<a$ the falling factorial is zero. These are finite polynomial, or Charlier-type, signed corrections; $H_0=1$ and $H_1=\kappa-r-2s$.

\begin{corollary}[All fixed order summable mass expansion]\label{spm:cor:massall}
For every fixed $J\geq0$,
\begin{equation}\label{spm:eq:massall}
 \sum_{r,s\geq0}\left|p_n(r,s)
  -p_{r,s}\sum_{j=0}^J H_j(r,s)n^{-j/2}\right|
 =O_J(n^{-(J+1)/2}).
\end{equation}
\end{corollary}
\begin{proof}
Multiply \eqref{spm:eq:Qpolynomials} by the limiting Poisson pgf. Coefficient extraction of $w^a\zeta^b$ times that pgf gives $p_{r,s}(r)_a(s)_b/(\lambda^a\mu^b)$, yielding \eqref{spm:eq:Charlier}. At every fixed order the normalized remainder is $O_J(n^{-(J+1)/2})$ on a fixed radius-$\rho$ polydisc. The same geometric Cauchy sum used for \eqref{spm:eq:masscorrection} proves \eqref{spm:eq:massall}.
\end{proof}

\section{Elementary growth and controlled inversion}\label{sec:inverse}
Combining \eqref{eq:logI} and \eqref{eq:unmarked}, with $L=\log2$, gives
\begin{equation}\label{eq:logA}
 \log A_n=\frac n2\log\frac{n}{eL^2}+\sqrt n+c
                 +\frac b{\sqrt n}+\frac d n+O(n^{-3/2}),
\end{equation}
where
\begin{equation}\label{eq:cbd}
 c=\frac{L^2-1}{4}-\log(2\sqrt2 L),\quad
 b=\frac7{24}-\frac{L^2}{6},\quad
 d=-\frac7{48}-\frac{L^2}{4}+\frac{L^4}{24}.
\end{equation}
Thus the elementary leading equivalent is
\begin{equation}\label{eq:elementary}
 A_n\sim\frac{e^{(L^2-1)/4}}{2\sqrt2 L}
       n^{n/2}(eL^2)^{-n/2}e^{\sqrt n}.
\end{equation}
The all-order algorithm and the involution saddle provide every further logarithmic coefficient by formal logarithm at each fixed order, with the already proved remainder.

\subsection{A Lambert W seed and its corrections}
For $X\to\infty$, let $y=\log X$, and define
\begin{equation}\label{eq:Wseed}
 w=W\left(\frac{2y}{eL^2}\right),\qquad N=\frac{2y}{w},\qquad Q=w+1,
\end{equation}
where $W$ is the positive real branch of the Lambert function. Then
$F_0(N)=y$ for $F_0(x)=(x/2)\log(x/(eL^2))$. Set
\begin{align}
 \alpha&=-\frac2Q,\label{eq:alpha}\\
 \beta&=-\frac{2c}{Q}+\frac2{Q^2}-\frac2{Q^3},\label{eq:beta}\\
 \gamma&=-\frac{2b}{Q}
 +c\left(\frac2{Q^2}-\frac4{Q^3}\right)
 -\frac1{Q^3}+\frac{14}{3Q^4}-\frac4{Q^5}.\label{eq:gamma}
\end{align}

\begin{proposition}[Controlled smooth inverse]\label{prop:inverse}
Let $\nu(X)$ be the eventual unique real solution of
\[
 \frac x2\log\frac{x}{eL^2}+\sqrt x+c+bx^{-1/2}+dx^{-1}=\log X.
\]
Then
\begin{equation}\label{eq:inverse}
 \nu(X)=N+\alpha\sqrt N+\beta+\frac\gamma{\sqrt N}
                     +O\left(\frac1{NQ}\right).
\end{equation}
This is the inverse of the specified smooth approximation, not a definition of a noninteger combinatorial count.
\end{proposition}
\begin{proof}
The left side of the defining equation has positive derivative for all sufficiently large $x$ and tends to infinity. At $N$,
\[
 F_0'(N)=Q/2,\qquad F_0''(N)=1/(2N),\qquad
 F_0'''(N)=-1/(2N^2).
\]
Substitute $x=N+\alpha\sqrt N+\beta+\gamma/\sqrt N$, expand $F_0$, $\sqrt x$ and $x^{-1/2}$, and cancel successively the coefficients of $\sqrt N$, $1$, and $N^{-1/2}$. This gives \eqref{eq:alpha}--\eqref{eq:gamma}. The remaining residual is $O(N^{-1})$; division by a derivative comparable to $Q$ gives the stated error. In particular the $d/x$ term first changes the root at the retained error scale.
\end{proof}

The same procedure constructs all fixed orders. Given enough terms in \eqref{eq:logA}, the resulting smooth inverse has, for every fixed $J\geq0$, the form
\begin{equation}\label{eq:inverseall}
 N+\sum_{j=0}^J N^{(1-j)/2}R_j(Q)
       +O_J\left(\frac{N^{-J/2}}{Q}\right),
\end{equation}
where $R_0=\alpha$, $R_1=\beta$, $R_2=\gamma$. Each $R_j$ is rational in $Q$, with possible poles only at $Q=0$, and depends on finitely many logarithmic coefficients of $A_n$. Inductively substitute the previous terms, compute the first remaining coefficient, and cancel it by division by $F_0'(N)=Q/2$. Taylor expansion of $\sqrt x$ and inverse half-powers introduces only rational functions of $Q$. The mean value theorem then gives the error at that fixed order. The smooth root is always defined using enough logarithmic terms for the asserted accuracy.

\subsection{Why integer thresholds need a rounding qualification}
Define the eventual threshold
\[
 m(X)=\min\{n\geq1:A_n\geq X\}.
\]
The sequence is strictly increasing for $n\geq1$. Adding $1$ to the first diagonal entry is an injection from $\mathcal A_n$ into $\mathcal A_{n+1}$. The identity matrix of dimension $n+1$ is outside its image, since a packed matrix of total sum $n$ has dimension at most $n$.

A rounding-safe shorthand for \eqref{eq:inverse} is
\begin{equation}\label{eq:rounding}
 m(X)=\ceil{N+\alpha\sqrt N+\beta+\gamma/\sqrt N
                    +O(1/(NQ))}.
\end{equation}
Its meaning is that a perturbation of the indicated size exists \emph{before} the ceiling is taken. It does not claim that the ceiling of the bare displayed truncation is always correct. The target $X$ can lie arbitrarily close to a sequence value.

To justify the shorthand directly, apply \eqref{eq:logA} at the consecutive integers $m(X)-1,m(X)$. The smooth root lies between
$m(X)-1-O(n^{-3/2}/\log n)$ and $m(X)+O(n^{-3/2}/\log n)$, with $n\asymp m(X)$. A perturbation on that scale makes its ceiling equal to $m(X)$. Combining with \eqref{eq:inverse} gives \eqref{eq:rounding}. In particular the real approximation determines $m(X)$ up to $O(1)$; away from the integer-boundary error band it determines the ceiling.

\begin{theorem}[Two-ceiling threshold bounds]\label{thm:ceilings}
For every fixed $J\geq0$, write the all-order logarithmic approximation as
\[
 F_J(x)=\frac x2\log\frac{x}{eL^2}+\sqrt x+c
                 +\sum_{j=1}^J b_jx^{-j/2},\qquad p=(J+1)/2.
\]
There exist $B_J>0$ and $n_J$ such that
\begin{equation}\label{eq:logbounds}
 F_J(n)-B_Jn^{-p}\leq\log A_n\leq F_J(n)+B_Jn^{-p}
 \qquad(n\geq n_J).
\end{equation}
Let $\ell_J(y)$ and $r_J(y)$ be the eventual roots of
\begin{equation}\label{eq:bracketroots}
 F_J(\ell_J)+B_J\ell_J^{-p}=y,\qquad
 F_J(r_J)-B_Jr_J^{-p}=y.
\end{equation}
For sufficiently large $X$,
\begin{equation}\label{eq:ceilings}
 \ceil{\ell_J(\log X)}\leq m(X)\leq\ceil{r_J(\log X)}.
\end{equation}
If $x_J$ solves $F_J(x_J)=\log X$, then
\begin{equation}\label{eq:rootwidth}
 r_J(\log X)-\ell_J(\log X)
 =O_J\left(x_J^{-(J+1)/2}/\log x_J\right).
\end{equation}
\end{theorem}
\begin{proof}
The constants in \eqref{eq:logbounds} exist by the proved fixed-order remainder. Both functions $H_\pm(x)=F_J(x)\pm B_Jx^{-p}$ are eventually strictly increasing. Their roots satisfy $\ell_J\leq r_J$. Every sufficiently large integer $n<\ell_J$ has
$\log A_n\leq H_+(n)<\log X$, while $n=\ceil{r_J}$ satisfies
$\log A_n\geq H_-(n)\geq\log X$. For $X$ sufficiently large all smaller exceptional indices are also below the threshold. This proves \eqref{eq:ceilings}. The derivative is comparable to $\log x_J$, so the mean value theorem applied to the perturbations $\pm B_Jx^{-p}$ proves \eqref{eq:rootwidth}.
\end{proof}

When the two ceilings agree, their common value is the exact threshold. Otherwise they isolate the ambiguity caused by an integer boundary. This theorem proves the existence of the remainder constants and starting indices, but supplies no numerical values for them. It is an asymptotic bound, not a certified finite-input inverse program. Producing such a program requires effective versions of the global and local remainder estimates.

\section{Exact reproduction and its limits}\label{sec:computation}
The package separates exact finite checks, symbolic coefficient identities, noncertified floating diagnostics, and deterministic document reproduction. None of these checks replaces the analytic tail and remainder proofs.

\subsection{Independent counting algorithms}
For the unmarked generating function $F_k(z)=(1-z)^{-k}(1-z^2)^{-(k^2-k)/2}$,
\[
 \frac{F_k'(z)}{F_k(z)}=\frac{k+k^2z}{1-z^2}.
\]
Writing $d_n(k)=n!f_n(k,1)$ gives the exact polynomial recurrence
\begin{equation}\label{eq:exactrecurrence}
 d_0(k)=1,\quad d_1(k)=k,\quad
 d_{n+1}(k)=k d_n(k)+n(k^2+n-1)d_{n-1}(k).
\end{equation}
If $P_j(u)=(1+u)^{-1}\sum_{k\geq0}q^kk^j$, then
\[
 P_0(u)=1,\qquad
 P_m(u)=u\sum_{j=1}^m\binom mj P_{m-j}(u).
\]
Thus the geometric transform is the finite exact sum
\begin{equation}\label{eq:finite-transform}
 A_n(u,v)=\sum_{j=0}^n[k^j]f_n(k,v)P_j(u).
\end{equation}
The implementation retains the more general marked rational logarithmic derivative
\[
 (1-vz)(1-z^2)F_k'(z)
 =\bigl(kv+(k^2-k)z-vk^2z^2\bigr)F_k(z).
\]
For rational $v=V/D$, it stores $D^nn!f_n(k,V/D)$ as an integer polynomial; exact fractions are introduced only for the final normalization. All terms are finite.

A second algorithm does not use that recurrence. At each integer auxiliary dimension it counts diagonal and off-diagonal occupancies by stars and bars:
\begin{equation}\label{eq:starsbars}
 f_n(k,v)=\sum_{\ell+2j=n}
 \binom{k+\ell-1}{\ell}
 \binom{\binom k2+j-1}{j}v^\ell,
\end{equation}
where a zero number of cells contributes $1$ for zero occupancy and $0$ otherwise. Zero-line deletion then gives
\begin{equation}\label{eq:deletion}
 A_n(u,v)=\sum_{d=0}^n u^d\sum_{k=0}^d
                   (-1)^{d-k}\binom dk f_n(k,v).
\end{equation}
Replacing the two stars-and-bars factors by $\binom{k}{\ell}$ and $\binom{\binom k2}{j}$ gives the independent binary algorithm. Collision markers are checked independently by expanding the two numerator factors in \eqref{eq:collision-fn} before applying \eqref{eq:deletion}.

A third small-case algorithm recursively enumerates upper-triangular entry assignments, tracks occupied lines, and keeps the full tuple $(D,T,R,S)$. It neither uses the geometric transform nor the polynomial recurrence. This checks model conventions, trace parity, and both meanings of collision counts.

\subsection{The bounded exact verification scope}
The verifier's main workload has a hard cap of $n=640$. More expensive independent algorithms have smaller explicit caps. The fixed checks include:
\begin{itemize}
\item the published A138178 prefix through $n=24$, and recurrence versus independent zero-line deletion through $n=32$;
\item four rational dimension/trace marker pairs through $n=16$;
\item the A135588 binary prefix through $n=12$, and its equality with zero collision markers through that index;
\item independent exact collision factorial moments through $n=12$;
\item direct matrix enumeration and joint-marker checks through $n=6$;
\item involution recurrence versus its explicit factorial sum at four trace markers through the requested main index.
\end{itemize}
The exact receipt records larger integers with a documented signed, length-prefixed binary hash encoding, not an unrestricted decimal conversion. The cap is an implementation boundary, not a range restriction on the mathematical theorems and not a promise that arbitrary fixed-order coefficient computation is cheap.

The symbolic verifier compares the coupled Gaussian calculation of $C_1,C_2$ with a separate factorial-degree-moment expression. It also checks the collision and binary first corrections, the involution logarithmic coefficients from its recurrence, and a rational enclosure proving $1<L+L^2<2$, as used in the sharp total-variation constant. These are exact algebraic checks. They do not constitute a machine-checked proof of the asymptotic estimates.

\subsection{Floating diagnostics}
A separate program evaluates leading ratios, corrected residuals, and collision moments from exact rational data at high decimal working precision. The points $(u,v)=(1,1),(2,2),(1,1/2)$ provide different cancellation tests. The latter two are diagnostics at those fixed points, not an enlargement of the specifically stated complex neighborhood in Theorem~\ref{thm:main}. Decimal outputs are explicitly labeled noncertified: they are not interval enclosures, effective error bounds, or evidence of a proven numerical onset. No claim is made that the entire external OEIS b-file was independently compared.

\subsection{Rebuilding from an immutable package}
The TeX source and reproducible Python verifier are included with the PDF. The detailed commands and pinned dependency versions are in the package's README and computation notes. The build verifies a SHA-256 manifest before running and verifies that source bytes have not changed afterward. Output is written to a fresh external directory; packaged source paths are not overwritten.

Public bounded inputs are validated before expensive work, with explicit exceptions rather than assertions that disappear under Python optimization. Every Python subprocess is invoked with \texttt{-B}, independently of environment settings; guard tests also run in ordinary and \texttt{-O} modes. TeX runs with \texttt{-no-shell-escape}. The ZIP has sorted members and fixed metadata. The replay test extracts the actual ZIP twice, rebuilds it under ordinary and optimized Python, and compares output bytes and ZIP metadata as well as source immutability. This establishes reproducibility for the recorded toolchain, not for arbitrary future TeX, compression, or symbolic-algebra versions.

The proofs and the exact finite computations serve different purposes. The former establish all fixed orders and their uniformity; the latter catch algebra, indexing, model, and packaging errors in a transparent bounded experiment.

\section{Further questions}\label{spm:sec:outlook}
The marked expansion settles all fixed algebraic orders in the stated parameter domain, but it leaves several concrete problems.

\paragraph{Effective constants and certified thresholds.}
The proof separates a nonreal-pole bound, a small-dimension envelope, a gamma tail, an angular tail, and an integrable Taylor remainder. Making each constant explicit would turn Theorem~\ref{spm:thm:ceilings} into a finite-input threshold certificate. The main challenge is to avoid constants so large that the result is unusable. A floating residual table alone cannot certify an onset or a remainder.

\paragraph{The trace parity transition.}
The theorem assumes $v$ near $1$, where the opposite involution saddle is smaller by $e^{-c\sqrt n}$. When $v$ approaches zero on the scale $n^{-1/2}$, that suppression disappears. A uniform two-saddle treatment should describe the transition between even and odd total weights while maintaining absolute control of the gamma and pole errors. The zero collision markers treated here do not cross this trace-marker boundary.

\paragraph{Lattice local limits.}
The exact relation $T_n\equiv n\pmod2$ suggests a joint local limit on the correct sublattice, coupled to the dimension on its different scale. Such a result requires characteristic-function estimates away from the small complex neighborhood used for cumulants. A local Edgeworth expansion would also clarify how the limiting independence is modified at the first lattice-sensitive order.

\paragraph{Large deviations and changing marker domains.}
The present compact domains are fixed independently of $n$. Allowing $u$ to approach zero or infinity changes the auxiliary-dimension saddle, and allowing collision markers to grow with $n$ can favor an atypically large population of repeated cells. Determining the boundaries of these regimes requires a new balance of the generating-function terms, rather than substitution into a fixed-compact error estimate.

\paragraph{Higher collision corrections and positive approximations.}
Corollary~\ref{spm:cor:massall} supplies every fixed signed polynomial correction in $\ell^1$. It would be useful to organize these polynomials efficiently and determine when a positive compound-Poisson or mixed-Poisson approximation can match more than the first correction. The adjusted independent Poisson law in \eqref{spm:eq:adjustedmeans} achieves order $n^{-1}$ without claiming optimality. Counts of cells of entry at least $3$, extra units rather than repeated cells, and the maximum entry invite related but distinct markings.

\paragraph{Exponentially small sectors and high orders.}
The proof identifies suppressed opposite-saddle and nonreal-pole effects only through bounds. A finite algebraic truncation has an error larger than those effects, so it does not by itself isolate their amplitudes. Separately controlled sector definitions, high-order coefficient growth, and an optimal-truncation analysis would be needed for a transseries theory. None is asserted by the present all-fixed-order expansion.

\paragraph{Other matrix conventions.}
Unlabeled symmetric matrices, prescribed dimension, fixed trace, and decreasing-margin variants require their own exact transform and source comparison. Even when involution numbers reappear, the normalization and the collision rates can change. A systematic classification of which conventions preserve this coupled gamma/involution mechanism would be more informative than transferring formulas between superficially similar sequences.

\begin{writenote}[Added 7 October 2026, batch~113 (reciprocal note): a decreasing-margin variant]
One of the variants named in this paragraph is now treated in the
collection's report \file{a104779-total-kostka-sums} (bundle Report~237,
batch~113; its symbols carry the mark ``tks'' here). It counts the
symmetric nonnegative integer matrices of total sum $n$ without zero rows
whose row sums are weakly decreasing, with positions distinguished as here:
OEIS A104779, the sum $a^{\rm tks}_n$ of all entries of the Kostka matrix of
order~$n$. These matrices are among those counted by $A_n$. That report
uses no geometric transform and no ordered Bell pole: it starts from
Schwob's cycle-index identity, and a positive global support bound with
analytic coefficient stabilization at every fixed support gives, for every
fixed $S\ge0$ (its Theorem~1.1),
\[
 a^{\rm tks}_n=C_{\rm tks}\sum_{s=0}^{S}H^{\rm tks}_sI_{n-s}
 +O_S\bigl(I_nn^{-(S+1)/2}\bigr),\qquad
 C_{\rm tks}=\prod_{j\ge2}\Bigl(1-\frac1{j!}\Bigr)^{-1},
\]
with $I_n=I_n(1)$, $H^{\rm tks}_0=1$ and $H^{\rm tks}_1=H^{\rm tks}_2=0$.
The involution numbers reappear and the normalization changes, as this
paragraph anticipates: with Corollary~\ref{spm:cor:unmarked},
$a^{\rm tks}_n/A_n\sim2C_{\rm tks}e^{-L^2/4}L^{n+1}$, so the matrices with
decreasing margins are an exponentially small fraction of those counted
here. No result of this report is used there, and that report treats one
convention, not the classification asked for here, so the question stands.
\end{writenote}

The central analytic point is reusable: control the entire coefficient polynomial with a positive envelope, retain both local normalizers, and prove compact-marker estimates on the domain actually needed for the probabilistic conclusion. Here that combination yields a complete fixed-order theory for the count, two fluctuation scales, a summable collision law, and an inverse with its integer uncertainty kept explicit.

\section{Further questions and research}\label{spm:sec:further}
\begin{writenote}[Added 6 October 2026, batch~108]
Following Vladimir Reshetnikov's standing rule of 4~October 2026, every
claim of the source, or of the write, that is not proved in full is stated
here as an open question, with its source, the argument offered and what
is missing; the source's own seven questions are those of
Section~\ref{spm:sec:outlook}, which stay as printed. The intake and the
write found no false statement in the source and narrowed none of its
claims; the write supplied one step that Section~\ref{spm:sec:saddles}
leaves implicit (the note at its end). The source imports no published
theorem, so there is no unchecked input.
\end{writenote}
\begin{enumerate}
\item\label{spm:q:third} \textbf{The third coefficient} (added by the
 write; Remark~\ref{spm:rem:third}). The value
 $C_3(L,1)=-L^2(50L^4+324L^2-1755)/6480$ comes from one SymPy evaluation
 of~\eqref{spm:eq:algorithm} by the write, which reproduces $C_1$ and $C_2$
 and is consistent with the exact counts through $n=500$ (table there),
 which certify nothing. \emph{Missing:} a
 second, independent derivation, for instance by the factorial-degree
 moments that the delivered \file{code/symbolic_coefficients.py} uses for
 $C_1$ and $C_2$, and the marked coefficient $C_3(L,v)$ away from $v=1$.
 The source computes the coefficients only through $C_2$ and claims
 nothing about $C_3$. \textup{(Added after the independent check of
 6~October 2026: Richardson extrapolation of exact counts to $n=2000$, and
 of marked counts at $u=\tfrac45,\tfrac54$ to $n=1500$, agrees with the
 formula for $C_3(L,1)$ to eight to ten digits at three values of $L$ (the
 note after Remark~\ref{spm:rem:third}); the derivation is still missing.)}
\item\label{spm:q:priority} \textbf{The literature boundary} (source:
 Section~\ref{spm:sec:model} and \file{SOURCES.md}). The source reports a
 bounded search that found no earlier asymptotic for A138178 and credits
 every prior result it found (Jovovic's transform,
 Cameron--Prellberg--Stark's binary equivalent and collision precedent);
 the write confirmed that neither OEIS entry states an asymptotic and that
 the two credited statements are as quoted (Remark~\ref{spm:rem:oeis}).
 \emph{Open:} whether the related matrix-composition and
 Burge-polynomial literature \cite{matrixcomp,caylerian}, which neither the
 write nor this report relies on, or work on normal semistandard tableaux,
 contains an asymptotic for the nonnegative symmetric count or its marked
 versions. The source claims no priority; what is missing is a search
 beyond the bounded one, and the write did not read those two papers.
\end{enumerate}

\begin{thebibliography}{9}
\addcontentsline{toc}{section}{References}
\bibitem{oeis138178}
The OEIS Foundation Inc., \emph{A138178: Symmetric nonnegative integer matrices without zero lines},
\url{https://oeis.org/A138178}. Definition, displayed terms, normal-tableau interpretation, and geometric generating-function formula credited to V. Jovovic (2009). Checked 5 October 2026.

\bibitem{oeis135588}
The OEIS Foundation Inc., \emph{A135588: Symmetric binary matrices without zero lines},
\url{https://oeis.org/A135588}. Binary companion sequence. Checked in the source review, 5 October 2026.

\bibitem{cps}
P. J. Cameron, T. Prellberg, and D. Stark,
\emph{Asymptotics for incidence matrix classes},
Electronic Journal of Combinatorics \textbf{13} (2006), R85.
\href{https://doi.org/10.37236/1111}{doi:10.37236/1111};
\url{https://arxiv.org/abs/math/0510155}.
The 4 April 2006 version was inspected. Proposition 4.2 proves the leading symmetric binary equivalent; Section 2 contains the nonsymmetric collision-Poisson precedent.

\bibitem{matrixcomp}
E. Munarini, M. Poneti, and S. Rinaldi,
\emph{Matrix compositions},
Journal of Integer Sequences \textbf{12} (2009), Article 09.4.8.
\url{https://cs.uwaterloo.ca/journals/JIS/VOL12/Rinaldi/rinaldi.pdf}.
A related ordered rectangular-matrix model; it is not the transpose-symmetric count studied here.

\bibitem{caylerian}
G. Cerbai and A. Claesson,
\emph{Enumerative aspects of Caylerian polynomials},
\url{https://arxiv.org/abs/2411.08426};
author version of 30 July 2025,
\url{https://akc.is/papers/047-Enumerative-aspects-of-Cay-polys.pdf}.
Related exact packed-matrix and Burge-polynomial framework.
\end{thebibliography}

\end{document}
